% =====================================================================
% WINE'26 submission draft -- Phase 1 restructure
% NOTE: working draft uses `article`; switch to \documentclass{llncs}
% (Springer LNCS) before submission. All structure is LNCS-compatible.
% Markers: [TODO: ...] = needs author action; [VERIFY] = proof written
% by drafting pass, must be independently checked before submission.
% =====================================================================
\documentclass[11pt]{article}
\usepackage[margin=1.1in]{geometry}
\usepackage{amsmath,amssymb,amsthm,mathtools}
\usepackage{booktabs}
\usepackage{enumitem}
\usepackage[colorlinks=true,linkcolor=blue,citecolor=blue,urlcolor=blue]{hyperref}
\usepackage{xcolor}

\newcommand{\todo}[1]{\textcolor{red}{[TODO: #1]}}
\newcommand{\verifyme}{\textcolor{orange}{[VERIFY]}}

\newtheorem{theorem}{Theorem}
\newtheorem{proposition}{Proposition}
\newtheorem{lemma}{Lemma}
\newtheorem{corollary}{Corollary}
\newtheorem{definition}{Definition}
\theoremstyle{remark}
\newtheorem{remark}{Remark}
\newtheorem{openproblem}{Open Problem}

\newcommand{\OPT}{\mathrm{OPT}}
\newcommand{\PW}{\mathrm{PW}}
\newcommand{\FAIR}{\mathrm{FAIR}}
\newcommand{\PoF}{\mathrm{PoF}}
\newcommand{\E}{\mathbb{E}}
\newcommand{\R}{\mathbb{R}}
\DeclareMathOperator*{\argmax}{arg\,max}
\DeclareMathOperator*{\argmin}{arg\,min}

\title{The Price of Fairness, Complexity, and Incentives\\
in No-Numeraire Combinatorial Auctions}
\author{Anonymous submission}
\date{}

\begin{document}
\maketitle

\begin{abstract}
Blockchain trade-intent auctions---most prominently CoW Protocol's deployed
CIP-67 \emph{fair combinatorial auction}---allocate batches of orders to
competing solvers under a structural constraint absent from classical
combinatorial auctions: \emph{there is no numeraire}. Value produced in one
trader's asset cannot be freely moved to another, so monetary transfers, the
basic tool of auction theory, are unavailable. Prior work characterizes the
$2\times 2$ economic equilibrium; the general-$n$ algorithmic theory is open.
We develop it along three axes, organized by a single invariant: the
complementarity factor $g$, the maximum ratio by which batching can amplify a
solver's standalone values. \textbf{(i) Price of fairness.} Under illiquidity
friction $\lambda=1-\theta$ (converting value across assets within a bid
incurs haircut $\lambda$), the worst-case welfare cost of the deployed
fairness filter is \emph{exactly}
\[
\PoF(\theta,g)=\min\!\left(g,\tfrac1\theta\right),
\]
tight for all $\theta$ and $g$, independent of the number of orders and of
prices; the no-numeraire ($\theta=0$) and full-numeraire ($\theta=1$)
endpoints give a dichotomy $\PoF\in\{g,1\}$. Batch size does not enter;
under a per-leg coverage floor $v_j\ge\alpha c_j$ the exact law is
$\min(g,\,\alpha+\frac{1-\alpha}{\theta})$---the price of fairness
prices the worst single-leg shortfall, discounted by liquidity. The
deployed filter is precisely the singleton-coalition core constraint; the
individual allocation always lies in the multiplicative $g$-core (tight),
and the no-numeraire core can be nonempty where the transferable core is
empty (a phenomenon of overlapping packages). Online, the price splits in two: against realized lifetime floors even
randomization pays exactly $g$, while with static public floors the
residual problem is one-way search over the spread---deterministic
$\Theta(\sqrt g)$, randomized $\Theta(\log g)$ on the single-contention
class; ratio-threshold greedy pays $\Omega(g)$, the swarm's online
incarnation. Manipulating the floors themselves is a priced game:
contamination has a closed-form steady-state bias with a sharp
participation threshold (only injection-cost floors deter; detection
merely tempers), the deployed rate limit acts exactly as a ramp tax on
attacker rents, and the bias feeds every floor-calibrated bound through
$g\mapsto g/(1-\mu)$.
Strategically, the
individual-only restriction of the first-price game has coarse-correlated
price of anarchy at most $e/(e-1)$---floors make standing bids coincide
with prices, a model-native smoothness ingredient---while batched bids
create a \emph{latent-batch} obstruction; the full-game price of anarchy
is conjectured $\Theta(g)$, with the lower bound $g$ realized at a weak
pure equilibrium. \textbf{(ii) Complexity.}
Winner determination is NP-hard already with the fairness filter vacuously
satisfied, yet fixed-parameter tractable in the number of directed token
pairs (small in deployment); the no-numeraire constraint reshapes the
economics but provably not the computational core. \textbf{(iii) Incentives.}
VCG is infeasible without a numeraire; best-bid fairness is incompatible with
strategyproofness already for a single order (under mild richness
conditions ruling out vacuous mechanisms), yielding a
feasibility/fairness/strategyproofness/efficiency trilemma. We quantify the
trilemma on the canonical single-batch screening core of the problem: the
price of strategyproofness there is $\Theta(\log g)$---a lower bound of
$\tfrac{g\ln g}{2(g-1)}$ against \emph{all} randomized dominant-strategy
mechanisms, proved by a one-sided (downward-IC) envelope argument forced by
the binding-bid model, matched on the core by a geometric posted-delivery
mechanism, with the deterministic optimum exactly $\sqrt g$. The lower
bound embeds into the full model, and the upper bound lifts to general
instances with single-minded solvers under public (historical)
benchmarks---the regime our impossibility results naturally single
out---via a universally truthful eligibility-posted packing mechanism; the
multi-minded case remains open. We calibrate $g$ and the directed-pair parameter on deployed
auction data, and discuss the online/repeated setting, where best-bid
fairness provably forces a historical benchmark.
\end{abstract}

% =====================================================================
\section{Introduction}
% =====================================================================

Trade-intent auctions are now core infrastructure of decentralized exchange:
traders submit \emph{intents} (``sell $X$ of asset $A$ for at least $Y$ of
asset $B$''), and professional intermediaries called \emph{solvers} compete
to execute batches of intents, returning surplus to traders. CoW Protocol's
deployed mechanism (CIP-67), which intermediates on the order of billions of
dollars monthly \todo{pull current volume figure + citation}, is a
\emph{fair combinatorial auction}: solvers submit individual-order bids and
batched bids; a batched bid is \emph{filtered} unless every order in it does
at least as well as that order's best standalone bid; the protocol then
selects the welfare-maximizing combination of surviving bids, subject to
uniform directed clearing prices (UDP).

Two features separate this object from the classical combinatorial-auction
literature. First, and structurally: \textbf{there is no numeraire}. Solvers
produce value \emph{in the traders' own assets}, the assets are illiquid, and
value produced in one trader's asset cannot be freely converted into
another's. Monetary transfers---the basic lubricant of mechanism design,
underlying VCG, core payments, and fairness-by-rebate---are simply
unavailable. Second, fairness here is \emph{endogenous}: the benchmark each
order must clear (its best standalone bid) is itself formed from the bids.

Canidio and Henneke~\cite{canidio2024} introduce the model, characterize the
$2\times2$ (two orders, two solvers) equilibrium, define best-bid fairness,
and observe that VCG is infeasible without a numeraire. The general-$n$
algorithmic theory---worst-case welfare guarantees, computational complexity,
and the cost of incentive compatibility---is left open. This paper develops
that theory.

\paragraph{One invariant.}
A single quantity organizes all of our results: the \emph{complementarity
factor} $g\ge 1$, the maximum ratio by which batching can amplify a solver's
standalone values (Definition~\ref{def:g}). We show that $g$ is
simultaneously: the exact price of fairness at the no-numeraire endpoint;
the exact class-level threshold of Walrasian/UDP existence (equilibria
exist for all instances iff $g=1$, by native no-numeraire arguments);
the worst-case gap that any online fair mechanism must pay (exactly $g$
against realized floors, even randomized; $\Theta(\sqrt g)$
deterministic and $\Theta(\log g)$ randomized once floors are static
and public); and the spread
parameter governing the $\Theta(\log g)$ price of strategyproofness. The
deployed mechanism is then best read as follows: \emph{whenever no
global UDP equilibrium exists, CIP-67 uses local per-bid UDP plus a fairness
filter to approximate an equilibrium that does not globally exist.}

\subsection{Our results}

\begin{table}[t]
\centering
\small
\begin{tabular}{@{}llll@{}}
\toprule
Axis & Result & Tightness & Where \\
\midrule
Fairness & $\PoF(\theta,g)=\min(g,1/\theta)$ & exact, $\forall\theta,g$ & Thm.~\ref{thm:pof} \\
 & Dichotomy: no numeraire $g$, numeraire $1$ & exact & Cor.~\ref{cor:dichotomy} \\
 & Batch size irrelevant; coverage law $\min(g,\alpha+\frac{1-\alpha}{\theta})$ & exact & Prop.~\ref{prop:dirrel}, Thm.~\ref{thm:coverage} \\
Core & Filter $=$ singleton-core; $g$-core nonempty, tight; NTU/TU separation & --- & Props.~\ref{prop:singletoncore}--\ref{prop:ntusep} \\
PoA & Individual-only CCE $\le e/(e-1)$; weak-PNE lower bound $g$; full game conj.\ $\Theta(g)$ & partial & Thm.~\ref{thm:indpoa}, Prop.~\ref{prop:poalower} \\
Online & Realized floors: rand.\ ratio exactly $g$; static floors: det.\ $\Theta(\sqrt g)$, rand.\ $\Theta(\log g)$; greedy $\Omega(g)$ & exact/class & Thms.~\ref{thm:onlinerand}--\ref{thm:onlinelb}, Prop.~\ref{prop:onlinegreedy} \\
Manipulation & Contamination: bias closed-form; threshold $\kappa_1\ge v\beta$ deters; rate limit $=$ ramp tax; floor-calibrated $g\mapsto g/(1-\mu)$ & --- & Props.~\ref{prop:contbias}--\ref{prop:contramp} \\
Structure & Walrasian/UDP exists for all instances iff $g=1$ (native proofs) & exact & Thm.~\ref{prop:walras} \\
Complexity & Winner determination NP-hard ($\forall g>1$) & --- & Thm.~\ref{thm:nphard} \\
 & FPT in \#directed pairs $P$ & --- & Prop.~\ref{prop:fpt} \\
Incentives & Best-bid fairness $\perp$ SP (already $n=1$) & --- & Thm.~\ref{thm:incompat} \\
 & Price of SP $=\Theta(\log g)$ on screening core; LB $\frac{g\ln g}{2(g-1)}$ embeds & up to $2e$ factor & Thm.~\ref{thm:posp} \\
 & Deterministic price of SP $=\sqrt g$ & exact & Thm.~\ref{thm:posp} \\
 & Lift: single-minded, public benchmark, $\Theta(1+\log g)$ & up to $\approx e$ factor & Thm.~\ref{thm:composition} \\
 & Multi-minded: independent prices $\Omega(g/\log^2 g)$; coupled menu $\Theta(1+\log g)$ & --- & Prop.~\ref{prop:indep}, Thm.~\ref{thm:coupled} \\
 & Multi-agent multi-minded: $O(N\log g)$; serial class $\Omega(\min(N,g))$ & --- & Thms.~\ref{thm:randprio},~\ref{thm:serialbarrier} \\
Online & Lifetime-benchmark online fairness pays exactly $g$ (det.) & exact & Prop.~\ref{prop:online} \\
\bottomrule
\end{tabular}
\caption{Summary of results. $g$ is the complementarity factor
(Def.~\ref{def:g}), $\theta=1-\lambda$ the liquidity (Def.~\ref{def:theta}),
$P$ the number of distinct directed token pairs.}
\label{tab:results}
\end{table}

\paragraph{(i) The exact price of fairness (Section~\ref{sec:pof}).}
Our first main result is a complete characterization of the welfare cost of
the fairness filter as a function of market liquidity. We model illiquidity
as a \emph{haircut}: within a winning bid, a solver may convert value across
the traders it serves, but transferring value $x$ out of one trader's asset
delivers only $(1-\lambda)x=\theta x$ in another's.
Theorem~\ref{thm:pof} shows
\[
\PoF(\theta,g) \;=\; \min\!\left(g,\ \tfrac1\theta\right),
\]
and this is tight: the bound is attained by explicit two-order instances,
and it is independent of $n$ and of the price vector. The two endpoints
recover a clean dichotomy (Corollary~\ref{cor:dichotomy}): with no numeraire
($\theta=0$) fairness costs exactly the full complementarity factor $g$;
with a numeraire ($\theta=1$) fairness is free. The proof of the upper bound
is a uniform \emph{per-bid repair} argument: each winning bid of the
unconstrained optimum is either repaired in place (converting surplus
traders' value into deficit traders' assets at the haircut) or dropped to
its individual allocation, and a deficit--surplus accounting shows the
delivered fraction is at least $\theta$ in either case. Notably, the
worst case is \emph{not} the maximally complementary bid ($\gamma=g$) one
would guess, but an extreme-asymmetry bid with gain exactly $1/\theta$
(Remark~\ref{rem:wrongguess}); and the continuous $\min(g,1/\theta)$ law is
a genuine consequence of modeling illiquidity as a haircut rather than as a
convertible fraction (Remark~\ref{rem:haircut}).

\paragraph{(ii) Where the new constraint does---and does not---bite
(Sections~\ref{sec:walras}--\ref{sec:complexity}).}
We prove a complete complexity dichotomy: the no-numeraire constraint
reshapes the \emph{economics} of the problem but provably creates no new
\emph{computational} hardness. Winner determination is NP-hard for every
fixed $g>1$, by a reduction from 3-dimensional matching that keeps the
fairness filter and the UDP constraint entirely vacuous
(Theorem~\ref{thm:nphard})---exactly the hardness profile of classical
combinatorial auctions, no more and no less. Conversely, the UDP constraint
itself is the source of tractability in deployment: under UDP all orders on
a directed token pair must clear in one bid, so the effective item count is
the number $P$ of distinct directed pairs, not $n$, and winner determination
is fixed-parameter tractable in $P$ (Proposition~\ref{prop:fpt}); we
calibrate $P$ on deployed data in Section~\ref{sec:empirics}. On the
structural side, UDP equilibria supporting efficient allocations exist
for \emph{every} instance iff $g=1$ (Theorem~\ref{prop:walras}), with both
directions proved natively in the no-numeraire model: total-value
subadditivity makes every batch deviation dominated under uniform pricing
(a one-line consequence of UDP), while for any $g>1$ individual
feasibility caps the supporting rates and a three-cycle of overlapping
batches refutes them. At the trader-optimal equilibrium the price system
\emph{is} the fairness-floor system, $w_j=c_j$. This is the formal sense
in which the deployed design approximates a nonexistent equilibrium. On
the strategic side, the individual-only restriction of first-price
competition under the filter is smooth---CCE price of anarchy at most
$e/(e-1)$ (Theorem~\ref{thm:indpoa}), since floors make standing bids
observable prices---while the full game is obstructed by latent batches
(Remark~\ref{rem:latentbatch}) and conjectured $\Theta(g)$.

\paragraph{(iii) The trilemma, quantified (Section~\ref{sec:incentives}).}
Without a numeraire VCG is infeasible (a winner may owe value in an asset it
does not produce; Proposition~\ref{prop:vcg}, due
to~\cite{canidio2024}). We prove the deeper obstruction: \emph{best-bid
fairness is incompatible with dominant-strategy incentive compatibility
already for a single order} (Theorem~\ref{thm:incompat})---the top producer,
by reporting truthfully, sets the very benchmark it must then clear, earning
zero; under any DSIC mechanism the effective benchmark collapses to the
runner-up level. This yields a four-way trilemma: feasibility, best-bid
fairness, strategyproofness, and efficiency are not jointly achievable, and
the three natural mechanisms each drop exactly one
(Remark~\ref{rem:trilemma}). Our second main result quantifies what
strategyproofness costs. Theorem~\ref{thm:posp} pins down the price of
strategyproofness---the best worst-case ratio of fair-optimal welfare to the
welfare of a feasible, fair, DSIC mechanism---on the canonical single-batch
screening core $\mathcal R_g$, to which we reduce: it is $\Theta(\log g)$
there, and the lower bound embeds into the full model since $\mathcal R_g$
is an instance subclass:
\begin{itemize}[leftmargin=1.5em]
\item a lower bound of $\frac{g\ln g}{2(g-1)}\ge\frac{\ln g}{2}$ against
\emph{all} randomized DSIC mechanisms, proved by reducing any such mechanism
to a menu over (win-probability, expected-delivery) pairs via a
\emph{one-sided} (downward-IC) envelope inequality---the binding-bid model
only permits under-reports, so the standard two-sided Myerson machinery is
unavailable---and integrating the resulting pointwise constraint against
the kernel $V^{-2}$;
\item a matching $e(\lceil\ln g\rceil+1)$ upper bound on the core by a
uniform geometric posted-delivery mechanism, which lifts to general
instances with single-minded solvers under public benchmarks with
guarantee $\FAIR/(e(\lceil\ln g\rceil+1)+1)$
(Theorem~\ref{thm:composition}), via a universally truthful
eligibility-posted packing mechanism; and
\item an exactly-$\sqrt g$ characterization of the deterministic subclass.
\end{itemize}
The lower bound is, to our knowledge, the first incentive-compatibility
impossibility proved \emph{inside} the no-numeraire model rather than
imported from digital-goods auctions: the screening variable is the
solver's deliverable value itself, and the fairness floor enters as a
hard delivery constraint rather than a revenue objective.

\paragraph{(iv) Deployment calibration and the online setting
(Sections~\ref{sec:empirics}--\ref{sec:online}).}
Because every bound in the paper is parameterized by $g$ and $P$, the theory
is falsifiable on deployment data. Section~\ref{sec:empirics} estimates the
empirical distributions of (lower bounds on) $g$ and of $P$ from CoW
Protocol auction data \todo{companion data import}, showing where deployed
markets sit on the $\PoF(\theta,g)$ surface and why the FPT bound of
Proposition~\ref{prop:fpt} explains observed tractability. Finally,
Section~\ref{sec:online} shows that in the repeated setting, under a
strong lifetime-scoped benchmark, adversarial online arrival forces any
deterministic best-bid-fair mechanism into individual-only allocation, at
competitive ratio exactly $g$
(Proposition~\ref{prop:online}); online fairness must therefore run against
a \emph{historical} benchmark, and we discuss the contamination-robustness
of the deployed quantile benchmark and the role of its downward rate limit,
deferring the full game-theoretic treatment as an open problem.

\subsection{Related work}\label{sec:related}

\paragraph{Trade-intent auctions and CIP-67.}
Canidio and Henneke~\cite{canidio2024} introduce fair combinatorial
auctions for trade intents, prove the $2\times2$ equilibrium
characterization, define best-bid fairness and the no-numeraire feasibility
constraint, and observe VCG infeasibility. We take their model as given and
develop the general-$n$ worst-case, complexity, and mechanism-design theory.
\todo{verify exact title/venue of arXiv:2408.12225 and update citation;
add follow-up literature since 2024.}

\paragraph{Combinatorial auctions.}
Complexity and approximation of winner determination, and the gap between
LP relaxations and integral allocations, are classical
\cite{nisan2006,lehmann2002,bikhchandani2002}; all of it presumes a
numeraire. Our Theorem~\ref{thm:nphard} shows the hardness core is
unchanged; the novelty is entirely on the welfare/incentive side.

\paragraph{Price of fairness.}
The price-of-fairness program of Bertsimas, Farias, and
Trichakis~\cite{bertsimas2011} and the fair-division literature
\cite{caragiannis2012} study welfare lost to fairness constraints in
divisible-resource and cake-cutting settings; the no-numeraire combinatorial
setting, where fairness interacts with asset-specific production, has not
been analyzed. The exact $\min(g,1/\theta)$ law and its haircut mechanism
appear to be new.

\paragraph{Posted prices and competitive auctions.}
Our $\Theta(\log g)$ result is structurally reminiscent of competitive
digital-goods auctions \cite{goldberg2006} and single-sample/posted-price
guarantees; the difference is that here the ``price'' is a \emph{delivery
obligation} constrained below by an endogenous fairness floor and above by
feasibility in a specific asset, and the lower bound must be proved over
mechanisms that allocate, not extract revenue. \todo{position against any
recent work on mechanism design without transfers / with burned money.}

\paragraph{Companion empirical work.}
A companion measurement paper \cite{companion} constructs and stress-tests
the historical fairness benchmark deployed on CoW post-CIP-67; the present
paper is its theoretical counterpart. The contributions are disjoint:
no result here appears there. \todo{anonymized cross-reference per WINE
double-blind policy.}

% =====================================================================
\section{Model}\label{sec:model}
% =====================================================================

\paragraph{Orders and traders.}
There are $n$ orders (traders) $[n]$. Trader $j$ wishes to acquire asset
$a_j$; its utility is the quantity of $a_j$ it receives. Assets have
common-knowledge notional prices $p_a>0$; delivering $x$ units of $a_j$ to
trader $j$ creates value $x\,p_{a_j}$. Each order specifies a directed token
pair (sell asset, buy asset).

\paragraph{Solvers and bids.}
A finite set of solvers intermediates. A \emph{bid} is a triple $(i,S,y)$:
solver $i$, order set $S\subseteq[n]$, and a production vector
$y=(y_j)_{j\in S}$, where $y_j$ is the quantity of $a_j$ solver $i$ delivers
to $j$ if it wins exactly $S$. Write $v_j:=y_j\,p_{a_j}$ for the produced
value of order $j$ and $V_i(S):=\sum_{j\in S}v_j$ for the bid value. Bids
with $|S|=1$ are \emph{individual}, otherwise \emph{batched}. Bids are
binding commitments: a winning solver must deliver its stated production
(settlement fails otherwise), so a solver never reports a production vector
it cannot produce; under-reporting is unrestricted.

\paragraph{No-numeraire feasibility.}
Production is asset-specific. Absent conversion, a winning solver on bid
$(i,S,y)$ delivers to each $j\in S$ some value $w_j\le v_j$ \emph{in asset
$a_j$}, and cannot move value across assets or traders. This is the
structural element absent from classical combinatorial auctions: there are
no monetary transfers.

\begin{definition}[Complementarity factor]\label{def:g}
The instance has complementarity factor $g\ge1$ if
$V_i(S)\le g\sum_{j\in S}V_i(\{j\})$ for every solver $i$ and every
$S\subseteq[n]$. Thus $g$ caps how much batching can amplify standalone
value; $g=1$ means additive (no complementarity).
\end{definition}

\begin{definition}[Fairness floor; fair allocation]\label{def:fair}
The \emph{floor} of order $j$ is $c_j:=\max_i V_i(\{j\})$, the value of its
best standalone bid; let $C:=\sum_j c_j$. An allocation is \emph{fair} if
every allocated order receives value $w_j\ge c_j$. Batched bids violating
this for some covered order are \emph{filtered}.
\end{definition}

\paragraph{Welfare.}
An allocation selects a set of winning bids covering disjoint order sets;
its welfare is $W=\sum_j w_j$ over allocated orders. Let $\OPT$ denote the
maximum welfare over feasible allocations ignoring fairness, and $\FAIR$ the
maximum over fair feasible allocations. The \emph{price of fairness} is
$\PoF:=\sup_{\text{instances}}\OPT/\FAIR$, with the supremum over a stated
instance class (we always make the class explicit).

\begin{definition}[Liquidity / illiquidity friction]\label{def:theta}
Within a single winning bid, a solver may convert value across the traders
it serves at haircut $\lambda\in[0,1]$: withdrawing value $x$ from one
trader's delivered asset yields $(1-\lambda)x$ of value in another's. Write
$\theta:=1-\lambda\in[0,1]$ for the \emph{liquidity}. $\theta=0$ is the pure
no-numeraire case; $\theta=1$ is a full numeraire. (Conversion is per-bid:
value cannot be moved between different winning bids.)
\end{definition}

\paragraph{Uniform directed prices (UDP).}
The deployed mechanism additionally requires uniform clearing prices per
directed token pair, which forces all orders on the same directed pair into
a single winning bid. UDP plays no role in Sections~\ref{sec:pof}
and~\ref{sec:incentives} (our constructions use distinct pairs, making it
vacuous) and reappears as the source of fixed-parameter tractability in
Section~\ref{sec:complexity}.

% =====================================================================
\section{The Exact Price of Fairness}\label{sec:pof}
% =====================================================================

This section proves the paper's first main theorem: a complete and tight
characterization of the welfare cost of the fairness filter as a function of
the liquidity $\theta$ and the complementarity factor $g$.

\begin{theorem}[Price of fairness, exact]\label{thm:pof}
Over instances with complementarity factor at most $g$ and liquidity
$\theta\in[0,1]$,
\[
\PoF(\theta,g)\;=\;\min\!\Big(g,\ \frac1\theta\Big)
\;=\;
\begin{cases}
g, & \theta\le 1/g,\\[2pt]
1/\theta, & \theta\ge 1/g,
\end{cases}
\]
with the convention $1/0=\infty$. The bound is independent of $n$ and of the
price vector, and is attained by two-order instances.
\end{theorem}

\begin{proof}
\emph{Upper bound, part 1: $\PoF\le g$.}
For any bid $(i,S,y)$, Definition~\ref{def:g} and $V_i(\{j\})\le c_j$ give
$V_i(S)\le g\sum_{j\in S}V_i(\{j\})\le g\sum_{j\in S}c_j$. Summing over the
winning bids of any feasible allocation, $\OPT\le gC$. The \emph{individual
allocation}---assigning each order $j$ to $\argmax_i V_i(\{j\})$---delivers
exactly $c_j$ to each $j$, is fair and feasible, and has welfare $C$; hence
$\FAIR\ge C$ and $\OPT/\FAIR\le g$.

\emph{Upper bound, part 2: $\PoF\le 1/\theta$ (per-bid repair).}
Assume $\theta>0$. Fix an optimal (unconstrained) allocation and consider
any one of its winning bids $(i,S,y)$, with value $V:=V_i(S)$ and floor mass
$C_S:=\sum_{j\in S}c_j$. If $C_S=0$ then every $c_j=0$, hence every solver's
standalone value for every $j\in S$ is $0$, and Definition~\ref{def:g}
forces $V=0$; such bids contribute nothing and may be ignored. So assume
$C_S>0$ and define the \emph{gain} $\gamma:=V/C_S$. Replacing the bid
by the individual allocation of $S$ shows (by optimality of the chosen
allocation among feasible ones, applied bid-locally) we may assume
$V\ge C_S$, and Definition~\ref{def:g} gives $V\le gC_S$; hence
$\gamma\in[1,g]$. We exhibit a fair completion of $S$ with welfare
$f\ge\theta V$:

\emph{Case $\gamma\le 1/\theta$: drop.} Replace the bid by the individual
allocation of $S$: fair, welfare $C_S=V/\gamma\ge\theta V$.

\emph{Case $\gamma>1/\theta$: repair in place.} Let
$d:=\sum_{j\in S}(c_j-v_j)^+$ be the total deficit and
$s:=\sum_{j\in S}(v_j-c_j)^+$ the total surplus, so $s-d=V-C_S$. Covering
the deficit requires withdrawing $d/\theta$ from surplus traders, feasible
iff $s\ge d/\theta$, i.e.\ iff $V-C_S\ge d\,\frac{1-\theta}{\theta}$. Since
$d\le C_S$ and $\gamma>1/\theta$ implies
$V-C_S>C_S\frac{1-\theta}{\theta}\ge d\,\frac{1-\theta}{\theta}$,
the repair is always feasible in this case. The repaired bid is fair and
delivers
\[
f \;=\; V-\frac{d}{\theta}(1-\theta)
\;\ge\; V-\frac{C_S(1-\theta)}{\theta}
\;=\; V\Big(1-\frac{1-\theta}{\theta\gamma}\Big)
\;\ge\; V\Big(1-\frac{1-\theta}{\theta\cdot(1/\theta)}\Big)
\;=\;\theta V,
\]
the last inequality because $1-\frac{1-\theta}{\theta\gamma}$ is increasing
in $\gamma$ and $\gamma\ge1/\theta$.

Since the optimum's winning bids cover disjoint order sets, applying the
better of drop/repair to each bid yields a fair feasible allocation of
welfare $\sum f\ge\theta\sum V=\theta\,\OPT$, so $\FAIR\ge\theta\,\OPT$ and
$\PoF\le1/\theta$.

\emph{Lower bound: a single extreme-asymmetry family attains
$\min(g,1/\theta)$.}
Set $\gamma^\ast:=\min(g,\,1/\theta)\in[1,g]$. Two orders on distinct
directed pairs (UDP vacuous). Order $1$ has no standalone bid ($c_1=0$);
order $2$ has a standalone bid of value $c_2=C>0$ from a rival solver. One
batched bid on $\{1,2\}$, from a solver with standalone capabilities
$V(\{1\})=0$ and $V(\{2\})=C$, produces total value $V=\gamma^\ast C$
\emph{entirely in trader $1$'s asset}: $v_1=\gamma^\ast C$, $v_2=0$. The
complementarity bound holds: $V=\gamma^\ast C\le g(0+C)$. Then
$\OPT=\gamma^\ast C$, using the batch.

For the batch to be fair, trader $2$ must receive at least $c_2=C$, which
requires withdrawing $C/\theta$ from trader $1$'s delivered asset (for
$\theta=0$ no value can be moved across traders at all, so the batch is
unfair whenever $c_2>v_2=0$ and the conventions $1/\theta=\infty$,
$C/\theta=\infty$ below are literal). If
$\theta<1/g$ then $\gamma^\ast=g$ and the available production
$\gamma^\ast C=gC<C/\theta$ is insufficient: the batch admits no fair
repair at all. If $\theta\ge1/g$ then $\gamma^\ast=1/\theta$ and the
required withdrawal $C/\theta$ equals the \emph{entire} production, so the
unique fair use of the batch delivers exactly $C$ (all to trader $2$,
trader $1$ getting $0\ge c_1$). In both regimes every fair allocation has
welfare exactly $C$ (the individual allocation also delivers $C$), so
$\FAIR=C$ and
\[
\frac{\OPT}{\FAIR}\;=\;\gamma^\ast\;=\;\min\!\Big(g,\frac1\theta\Big).
\qedhere
\]
\end{proof}

\begin{remark}[Scope of the lower-bound family]\label{rem:nodominance}
The construction uses a solver whose batched production vector does not
componentwise dominate its standalone productions (standalone it can
deliver $C$ to order $2$; batched, it delivers order $2$ nothing).
Definition~\ref{def:g} constrains only total values, so this is admissible;
it also reflects practice, where batched routing genuinely reallocates
which legs produce surplus. Under an additional componentwise
free-disposal axiom ($v_j(S)\ge V_i(\{j\})$ for $j\in S$) the family is
excluded and the tight constant may change; we flag this as a deliberate
modeling choice rather than an oversight.
\end{remark}

\begin{remark}[A correction to the natural construction]
\label{rem:concentration}
The natural ``concentrated production'' instance (symmetric floors
$c_1=c_2=v$, batch of value $2gv$ all in one asset) attains ratio $g$ only
for $\theta\le\frac{1}{2g-1}$: for $\theta\in(\frac{1}{2g-1},\frac1g]$ its
repair is feasible and strictly beats dropping, so it certifies strictly
less than $g$. The extreme-asymmetry family above is what makes the lower
bound tight throughout $\theta\le1/g$; asymmetry of the \emph{floors}, not
magnitude of the complementarity, is the binding obstruction.
\end{remark}

\begin{corollary}[Dichotomy]\label{cor:dichotomy}
With no numeraire ($\theta=0$), $\PoF=g$ exactly. With a full numeraire
($\theta=1$), $\PoF=1$: the unconstrained optimum can always be
redistributed to clear every floor, since $\OPT\ge C=\sum_j c_j$.
\end{corollary}

\begin{remark}[The natural guess is wrong]\label{rem:wrongguess}
One might guess the worst case concentrates at maximal gain $\gamma=g$,
yielding $\PoF=\frac{g\theta}{(g+1)\theta-1}$. It does not: the
extreme-asymmetry bid with $\gamma=1/\theta$ is strictly worse whenever
$\theta\in(1/g,1)$, with gap governed by $(1-\theta)(g\theta-1)\ge0$. The
binding instance sits at the \emph{repair/drop indifference point}, not at
maximal complementarity. Section~\ref{sec:coverage} pins down what the
law prices: batch size is irrelevant (Proposition~\ref{prop:dirrel}),
and under a per-leg coverage floor $\alpha$ the exact three-parameter
law is $\PoF(\theta,g,\alpha)=\min(g,\alpha+\frac{1-\alpha}{\theta})$
(Theorem~\ref{thm:coverage}).
\end{remark}

\begin{remark}[Haircut vs.\ convertible fraction]\label{rem:haircut}
If instead a fixed fraction $\theta$ of a bid's value were
\emph{frictionlessly} reallocatable, $\PoF$ would be a step function ($1$ if
$\theta\ge1/g$, else $g$). The continuous law $\min(g,1/\theta)$ is a
genuine consequence of modeling illiquidity as a haircut---i.e., of slippage
being proportional to the amount converted---which is the empirically
correct model for AMM-based conversion. \todo{one sentence + citation on
AMM slippage linearity for small trades.}
\end{remark}

% =====================================================================
\subsection{What the price of fairness prices: coverage, not size}
\label{sec:coverage}

Two refinements pin down which feature of batching the
$\min(g,1/\theta)$ law actually prices.

\begin{proposition}[Batch size is not the parameter]\label{prop:dirrel}
Let $\PoF(\theta,g,d)$ be the price of fairness over instances whose bids
cover at most $d$ orders. Then $\PoF(\theta,g,1)=1$ and
$\PoF(\theta,g,d)=\min(g,1/\theta)$ for every $d\ge2$: the tight family of
Theorem~\ref{thm:pof} uses two-order batches, and the upper bound is
size-free.
\end{proposition}

What the construction exploits is not size but a single starved leg. The
right third parameter is \emph{per-order coverage}: say an instance has
coverage $\alpha\in[0,1]$ if every bid satisfies $v_j\ge\alpha c_j$ for
every order it covers---no leg is worse than $\alpha$ times the
benchmark, a per-leg routing-quality floor that is directly auditable in
deployment.

\begin{theorem}[The coverage law]\label{thm:coverage}
Over instances with liquidity $\theta$, complementarity factor at most
$g$, and coverage $\alpha$,
\[
\PoF(\theta,g,\alpha)\;=\;\min\!\big(g,\ \Gamma(\theta,\alpha)\big),
\qquad
\Gamma(\theta,\alpha)\;:=\;\alpha+\frac{1-\alpha}{\theta}
\;=\;1+\frac{(1-\alpha)(1-\theta)}{\theta},
\]
with the convention $\Gamma(0,\alpha)=\infty$ for $\alpha<1$ and
$\Gamma(0,1)=1$; tight on two-order instances. The endpoints recover
Theorem~\ref{thm:pof} ($\alpha=0$) and triviality at $\alpha=1$ or
$\theta=1$.
\end{theorem}

\begin{proof}
\emph{Upper bound.} Refine the per-bid repair/drop argument of
Theorem~\ref{thm:pof}. Fix an OPT bid with $C_S>0$ (degenerate bids as
before), gain $\gamma=V/C_S\in[1,g]$. Coverage caps the deficit:
$D=\sum_{j\in S}(c_j-v_j)^+\le(1-\alpha)C_S$. Repair withdraws
$D/\theta$ from the surplus $s=V-C_S+D$, feasible iff
$V-C_S\ge D\frac{1-\theta}{\theta}$, which is implied by
$\gamma\ge\Gamma$ since $D\le(1-\alpha)C_S$ and
$\Gamma-1=\frac{(1-\alpha)(1-\theta)}{\theta}$; the repaired welfare is
\[
f\;=\;V-\frac{D(1-\theta)}{\theta}
\;\ge\;V\Big(1-\frac{(1-\alpha)(1-\theta)}{\theta\gamma}\Big)
\;\ge\;\frac{V}{\Gamma}
\qquad(\gamma\ge\Gamma),
\]
the last step by monotonicity in $\gamma$ with equality at
$\gamma=\Gamma$. For $\gamma\le\Gamma$, dropping gives
$f=C_S=V/\gamma\ge V/\Gamma$; and for all $\gamma\le g$ dropping gives
$f\ge V/g$. Hence every bid retains a $1/\min(g,\Gamma)$ fraction, and
summing over the disjoint OPT bids, $\FAIR\ge\OPT/\min(g,\Gamma)$.

\emph{Lower bound (coverage-saturated asymmetry).} Set
$\gamma^\ast:=\min(g,\Gamma)$. Two orders on distinct pairs: $c_1=0$,
$c_2=c>0$; one batched bid with $v_2=\alpha c$ (the coverage bound
saturated: maximal allowed deficit) and $v_1=(\gamma^\ast-\alpha)c$, so
$V=\gamma^\ast c$, $C_S=c$; the batching solver's standalone values are
$0$ and $c$, so Definition~\ref{def:g} holds since $\gamma^\ast\le g$.
Repairing order $2$ requires withdrawing $(1-\alpha)c/\theta$ from
$v_1$. If $\gamma^\ast=\Gamma$, then
$v_1=(\Gamma-\alpha)c=\frac{(1-\alpha)c}{\theta}$ \emph{exactly}: the
unique fair use of the batch delivers exactly $c$. If
$\gamma^\ast=g<\Gamma$, then $v_1<\frac{(1-\alpha)c}{\theta}$ and the
batch is irreparable. In both regimes $\FAIR=c$ and
$\OPT/\FAIR=\gamma^\ast$ ($\theta=0$ read with the conventions of
Theorem~\ref{thm:pof}).
\end{proof}

\begin{remark}[Reading the law]\label{rem:coveragereading}
The price of fairness prices neither batch size
(Proposition~\ref{prop:dirrel}) nor complementarity per se:
$\Gamma-1=(1-\alpha)\cdot\frac{1-\theta}{\theta}$ is the largest
relative per-leg deficit $(1-\alpha)$ multiplied by the conversion loss
per unit of delivered deficit, $(1-\theta)/\theta$. The endpoint
$\alpha=1$ is exactly the componentwise filter-passing \emph{part} of
assumption (A2) in Section~\ref{sec:composition}: in that regime
fairness is free at the welfare level, and Section~\ref{sec:composition}
then asks what strategyproofness costs on top of that fair-feasible
baseline. Deployed per-leg quality floors thus directly cap the welfare
cost of the fairness filter. \todo{Section~\ref{sec:empirics}: estimate
the empirical coverage distribution $\widehat\alpha$ alongside
$\widehat g$.}
\end{remark}

\section{One Invariant: \texorpdfstring{$g$}{g} as the Walrasian Gap}\label{sec:walras}
% =====================================================================

The constant $g$ in Theorem~\ref{thm:pof} is not an artifact of the fairness
analysis; it is the market's fundamental complementarity invariant. View the
setting as an exchange market: traders hold assets and want others, solvers
intermediate, and UDP prices are candidate uniform clearing prices.

\paragraph{UDP equilibrium without a numeraire.}
Uniform directed prices assign a rate to each directed token pair; under
the notional prices, the rates determine for each order $j$ a value
$w_j\ge0$ that \emph{any} solver serving $j$ must deliver to $j$, in asset
$a_j$---the single property of UDP used below is that $w_j$ depends only
on the order, not on the bid serving it. A solver's feasible action at
rates $w$ is to execute a disjoint portfolio of its bids; a batched bid
$(i,S,y)$ is \emph{executable} iff $v_j\ge w_j$ for every $j\in S$ (at
$\theta=0$ no value moves across assets), with profit
$\sum_{j\in S}(v_j-w_j)$, retained in the respective assets and compared
by notional value. A \emph{UDP equilibrium supporting the efficient
allocation} is a pair (rates, allocation) such that every order is served
and receives exactly $w_j\ge0$; every solver's executed portfolio
maximizes its profit over feasible portfolios at $w$; and the allocation
is \emph{production-efficient}: its produced value equals $\OPT$. In
this subsection an allocation's social surplus is thus its produced value,
which the obligations $w_j$ split between trader deliveries and solver
surplus; individual rationality and fairness are checked through the
$w_j$. We work in the order-price abstraction of UDP: a directed-rate
system induces order-level obligations $w_j$, and statements are for
rate-realizable obligation vectors.

\begin{lemma}[Batch domination under $g=1$]\label{lem:domination}
If the instance satisfies Definition~\ref{def:g} with $g=1$, then at any
rates $w$, every executable batched bid of every solver is weakly
dominated by the sub-portfolio of its positive-profit singletons:
\[
\sum_{j\in S}(v_j-w_j)\;=\;V_i(S)-\sum_{j\in S}w_j
\;\le\;\sum_{j\in S}V_i(\{j\})-\sum_{j\in S}w_j
\;\le\;\sum_{j\in S}\big(V_i(\{j\})-w_j\big)^+,
\]
where each positive-profit singleton is automatically executable
($V_i(\{j\})>w_j$); a non-executable batch is not a deviation at all.
Hence every solver has an individual-only best response. The single use
of UDP is that $w_j$ is the same whether $j$ is served in the batch or
individually.
\end{lemma}

\begin{theorem}[The Walrasian/UDP existence frontier is $g=1$]
\label{prop:walras}
The class of instances with complementarity factor at most $g$ admits a
UDP equilibrium supporting an efficient allocation \emph{for every
instance} iff $g=1$. Precisely: (a) every instance satisfying the
total-value bound of Definition~\ref{def:g} with $g=1$ (no per-order
additivity required) admits such an equilibrium, in which the obligations
are exactly the fairness floors, $w_j=c_j$, so the equilibrium is best-bid
fair; (b) for every $g>1$ there exists an instance admitting none.
(Individual $g>1$ instances may of course admit equilibria.) Throughout,
solvers have no capacity constraints (a solver may execute any disjoint
portfolio, consistent with Sections~\ref{sec:incentives}), and in (a) the
floor vector $w_j=c_j$ is assumed rate-realizable---automatic when orders
lie on distinct directed pairs, or when the public benchmark is itself
produced by a UDP reference outcome.
\end{theorem}

\begin{proof}
\emph{(a)} Set $w_j:=c_j$ and assign each order $j$ to a solver attaining
$V_i(\{j\})=c_j$. The winner of each order earns zero margin; any
losing singleton has $V_i(\{j\})-c_j\le0$; and by
Lemma~\ref{lem:domination} every batch deviation is dominated by its
positive-profit singletons, of which there are none. So no solver has a
profitable deviation. The allocation is production-efficient: with $g=1$
any allocation produces at most
$\sum_{(i,S)}V_i(S)\le\sum_{(i,S)}\sum_{j\in S}V_i(\{j\})\le
\sum_jc_j=C$, and this allocation produces exactly $C=\OPT$, all of it
delivered. Finally $w_j=c_j$ is the best-bid fairness floor, so the
equilibrium is fair. (The construction is the trader-optimal core point of
the underlying Shapley--Shubik assignment market \cite{shapley1971},
to which Lemma~\ref{lem:domination} reduces the economy.)

\emph{(b)} We argue natively; note that non-existence in the
\emph{with-transfers} relaxation does \emph{not} imply non-existence
here, since the no-numeraire constraint only removes deviations
(executability), making allocations easier to support---so an LP
integrality gap alone is insufficient evidence. Take the three-cycle: fix
$\delta\in(0,2(g-1)]$; orders $1,2,3$ on distinct pairs; a rival solver
with standalone value $1$ on each order ($c_j=1$); three solvers, each
holding one batched bid on a pair $\{j,j'\}$ with production vector
$(1+\delta/2,\,1+\delta/2)$ (value $2+\delta\le2g$) and standalone values
$1$ on its two orders. The efficient allocation executes one batch, say
$\{1,2\}$, plus order $3$ individually, with produced value $3+\delta$.
Suppose rates $w$ support it. Order $3$ is served individually by a
solver of standalone value $1$, so feasibility forces $w_3\le1$; likewise
$w_1,w_2\le1+\delta/2$ (and, being served, $w_1,w_2\le$ the serving
batch's per-order production). Since $w_2\le1+\delta/2$ and
$w_3\le1<1+\delta/2$, the bid of the solver holding batch $\{2,3\}$
is executable, and equilibrium requires it unprofitable relative to that
solver's individual alternatives:
\[
(2+\delta)-w_2-w_3\;\le\;(1-w_2)^++(1-w_3)^+ .
\]
If $w_2\le 1$ the right side is $(1-w_2)+(1-w_3)$ and the inequality
gives $\delta\le0$, a contradiction. If $w_2>1$, the right side is
$1-w_3$, giving $w_2\ge1+\delta$; but order $2$ is served by batch
$\{1,2\}$ whose per-order production is $1+\delta/2<1+\delta$,
contradicting feasibility. Either way no supporting rates exist.
\end{proof}

\begin{remark}[The with-transfers benchmark]\label{rem:botransfers}
In the with-transfers relaxation, the Bikhchandani--Ostroy criterion
\cite{bikhchandani2002} characterizes existence by LP integrality, and
$g=1$ implies integrality directly (shift any fractional weight from a
batch to its singletons: feasible and weakly improving). The three-cycle
exhibits the strict gap
$\tfrac32(2+\delta)>3+\delta$. We record two structural notes. First,
single-batch instances such as those of Theorem~\ref{thm:pof} have
integral LPs: \emph{overlap} among batches, not complementarity alone,
breaks existence. Second, in the no-numeraire UDP model the Bikhchandani--Ostroy
implication cannot be used in either direction: feasibility constraints
remove deviations (aiding existence), while the absence of cross-asset
transfers removes supporting prices---which is why both directions of
Theorem~\ref{prop:walras} required native proofs, and why we do not state
a per-instance LP criterion for the no-numeraire notion.
\end{remark}

\begin{remark}[One invariant, several faces]\label{rem:faces}
Three a-priori distinct quantities are governed by $g$: (i) the price of
fairness at $\theta=0$ (Theorem~\ref{thm:pof}, tight); (ii) the LP
integrality gap of welfare maximization, which is at most $g$ (the LP
optimum is $\le gC\le g\cdot$ integral OPT, by the proof of
Theorem~\ref{thm:pof}), the with-transfers existence criterion
(Remark~\ref{rem:botransfers}); (iii) the core deficit of the
efficient allocation: the floors $c_j$ are precisely the individual
rationality constraints, and $\OPT-C\le(g-1)C$. These are not equal
instance-by-instance; $g$ is the common worst-case governor, witnessed by
\emph{different} extremal obstructions: floor asymmetry for the price of
fairness (Remark~\ref{rem:concentration}) and overlapping packages for
Walrasian failure (Theorem~\ref{prop:walras}(b))---single-batch
instances have integral LPs. This reads the deployed design
as principled: \emph{whenever no global UDP equilibrium exists, CIP-67
selects a fair point via local per-bid UDP plus the filter---approximating
an equilibrium that does not globally exist.}

A fourth face---the strategic price of the first-price fair combinatorial
auction---is treated next: it provably decomposes into the filter's $g$
times a smooth constant for individual competition, and conjecturally so
in general.
\end{remark}

\subsection{The strategic face: smooth individuals, latent batches}
\label{sec:poa}

We study the first-price (pay-as-bid) fair combinatorial auction: solvers
submit binding bids (individual and batched, with committed deliveries
bounded by capability), batched bids failing any covered floor are
filtered, and a welfare-maximizing disjoint collection of surviving bids
is executed, each winner delivering its bid. All results are for
\emph{produced} welfare, denoted
$\PW(b):=\sum_{j\text{ won}}v_{\mathrm{winner}(j),j}$ to distinguish it
from the delivered welfare $W$ of the rest of the paper; with
$R(b):=\sum_{j\text{ won}}p_j(b)$ we have
$\sum_iu_i(b)=\PW(b)-R(b)$. Produced welfare is the standard auction
objective; the delivered-welfare price of anarchy is unbounded (literally
infinite if zero bids are allowed, otherwise via $\varepsilon$-bids by a
lone solver producing its full capability).

Two model features do real work. First, bids are binding, so
\emph{no-overbidding}---ordinarily an assumption in price-of-anarchy
analyses---is structural here. Second, in the \emph{individual-only} game
(all bids individual), three quantities coincide for every profile $b$
and order $j$:
\[
c_j(b)\;=\;\max\{\text{standing individual bids on }j\}\;=\;p_j(b),
\]
the floor, the best standing offer, and the delivered amount (welfare
maximization over individual bids decomposes per order; the maximal bid
wins and delivers itself). Standing offers \emph{are} prices: because the
fairness benchmark is computed from the bids, the competitive price is
observable in the profile. This identity is exactly the smoothness
ingredient whose absence for batch components obstructs the general case.

\begin{theorem}[Individual competition is smooth]\label{thm:indpoa}
Consider the individual-only game, assumed per-order decomposable: either
UDP is vacuous (orders on distinct directed pairs) or each directed pair
is aggregated into one atomic item; solvers have additive, no-capacity
individual production across these items, and zero-value optimal orders
are ignored. For every coarse correlated equilibrium,
\[
\E[\PW]\;\ge\;\Big(1-\frac1e\Big)\,\OPT_{\mathrm{ind}},
\qquad
\OPT_{\mathrm{ind}}:=\sum_j v^{\max}_j,
\]
so the CCE price of anarchy of the individual-only restriction is at most
$e/(e-1)\approx1.582$.
\end{theorem}

\begin{proof}
For order $j$ let $i^\ast(j)$ attain $v^{\max}_j$ and let this solver
deviate, independently across its orders, to the random bid $W_j$ with
density $f(w)=1/(v^{\max}_j-w)$ on $[0,(1-1/e)v^{\max}_j]$ (the
Syrgkanis--Tardos first-price deviation \cite{syrgkanis2013}). Individual
bids are never filtered and bids on distinct orders never conflict, so
utilities add over orders, and the deviator wins $j$ whenever
$W_j$ exceeds $\beta_j:=\max\{\text{standing bids on }j\text{ in
}b_{-i}\}$:
\[
\E_W\big[u_{i^\ast}\text{ from }j\big]
=\int_{0}^{(1-1/e)v^{\max}_j}\!\!(v^{\max}_j-w)f(w)\,
\mathbf 1\{w>\beta_j\}\,dw
=\Big((1-\tfrac1e)v^{\max}_j-\beta_j\Big)^{\!+}
\ge(1-\tfrac1e)v^{\max}_j-p_j(b),
\]
using $\beta_j\le c_j(b)=p_j(b)$, the identity above---the penalty lands
under the original profile $b$ ($f$ integrates to
$\ln\frac{v}{v-(1-1/e)v}=1$). Randomized deviations are allowed at CCE
by linearity: fix the deviation distribution and take expectation over
its internal seed. Summing over orders and solvers,
$\sum_iu_i(\mathrm{dev}_i,b_{-i})\ge(1-1/e)\OPT_{\mathrm{ind}}-R(b)$.
At a CCE, $\E[\sum_iu_i(b)]\ge\E[\sum_iu_i(\mathrm{dev}_i,b_{-i})]$,
and $\sum_iu_i(b)=\PW(b)-R(b)$; the two $R(b)$ terms cancel.
\end{proof}

\begin{corollary}[Against the unrestricted benchmark]\label{cor:indpoa}
If the auction is restricted to individual bids, its CCE produced welfare
is at least $(1-1/e)\OPT_{\mathrm{ind}}$, hence at least
$(1-1/e)\,\OPT/g$ against the unrestricted production benchmark, since
$V_i(S)\le g\sum_{j\in S}V_i(\{j\})\le g\sum_{j\in S}v^{\max}_j$
summed over the optimal partition gives $\OPT\le g\,\OPT_{\mathrm{ind}}$.
This is a guarantee for the individual-only restriction; it is \emph{not}
a full-game CCE bound---extending it is precisely the latent-batch
obstruction below.
\end{corollary}

\begin{proposition}[The filter's $g$, strategically dressed]
\label{prop:poalower}
At $\theta=0$, measured against the \emph{unrestricted} produced
optimum, the full game admits a \emph{weak} pure Nash equilibrium with
$\PW=\OPT/g$: orders $1,2$, a rival with standalone capabilities
$(v,v)$, and a solver $A$ with standalone $(v,v)$ and batch capability
$(2gv,\,0)$. All solvers bidding their standalone capabilities truthfully
(batch unsubmitted) is a weak equilibrium: the batch is filtered for
every delivery split and every submission strategy of $A$, since it
cannot give order $2$ its floor $v$, which the rival's standing bid pins
regardless of $A$'s behavior; individual deviations earn $0$ either way
at the capability tie (binding bids forbid exceeding $v$), under any
fixed tie-breaking. The loss is entirely the fairness filter's, not
strategy's: against the \emph{fair} optimum the same profile is optimal
---this is Theorem~\ref{thm:pof}'s instance wearing a strategic hat.
\end{proposition}

\begin{remark}[The latent-batch obstruction, and a conjecture]
\label{rem:latentbatch}
For the full game, every deviation-based route hits the same wall: an
individual deviation bid $w_j$ can lose only to a bid that
\emph{executes under the deviation profile} and delivers at least $w_j$
(the filter forces this), while smoothness penalties must be charged
under the \emph{original} profile. A \emph{latent batch}---present in
$b_{-i}$, executing under the deviation but not under $b$---defeats every
charging scheme we tried, including profile-dependent deviations at pure
equilibria. The failure mode is informative: when a latent batch eats the
deviator, the deviation profile's produced welfare is high, so the
deviation is socially excellent but privately worthless---a coordination
failure invisible to unilateral-deviation arguments. We conjecture that
the full-game CCE price of anarchy for produced welfare is $\Theta(g)$,
with the natural target upper bound $g\cdot e/(e-1)$. What is proved: the
lower bound $g$ (Proposition~\ref{prop:poalower}) and the upper bound for
the individual-only restriction (Theorem~\ref{thm:indpoa},
Corollary~\ref{cor:indpoa}). Extending the upper bound to the full game
is exactly the latent-batch obstruction: any proof must control latent
batches, and any counterexample must sustain them in CCE.
\end{remark}

\subsection{The fairness hierarchy and the no-numeraire core}
\label{sec:hierarchy}

The fairness filter has a cooperative-game reading that places it in a
hierarchy and exhibits, for a third time, the pattern that the
no-numeraire constraint \emph{removes deviations}.

\paragraph{Blocking.}
A bid $(i,S,y)$ \emph{blocks} an allocation $w$ at factor $\kappa\ge1$ if
some feasible delivery from $y$ (componentwise, with $\theta$-haircut
conversion within the bid) gives every $j\in S$ strictly more than
$\kappa\,w_j$. The \emph{$\kappa$-core} is the set of allocations blocked
by no bid at factor $\kappa$; the $1$-core is the (NTU) core.

\begin{proposition}[The filter is the singleton core constraint]
\label{prop:singletoncore}
An allocation is best-bid fair iff it is blocked by no \emph{standalone}
bid at factor $1$: $w_j\ge c_j$ for all $j$ iff no singleton coalition
improves. The deployed filter is exactly the singleton-coalition core
constraint; the full core additionally forbids batch blocking.
\end{proposition}

\begin{proposition}[The $g$-core is never empty; $g$ is tight]
\label{prop:gcore}
(a) For every $\theta\in[0,1]$ and every instance, the individual
allocation $w_j=c_j$ lies in the $g$-core: blocking at factor $g$ would
deliver strictly more than $g\,c_j$ to every $j\in S$, for a total
exceeding $g\,C_S$, while any feasible delivery from a bid totals at most
$V_i(S)\le g\,C_S$ (haircuts only lose value); if all floors in $S$ are
zero, Definition~\ref{def:g} forces $V_i(S)=0$ and strict blocking is
impossible. (b) Tight: for every
$\kappa<g$ there is an instance whose individual allocation is not in the
$\kappa$-core---a batch with componentwise vector $(g'c,\,g'c)$ over
floors $(c,c)$, any $g'\in(\kappa,g]$. Hence $g$ is exactly the least
universal blocking-immunity factor of the individual allocation; at $g=1$
the individual allocation is in the core, consistent with
Theorem~\ref{prop:walras}, where it is the equilibrium allocation.
\end{proposition}

\begin{proposition}[TU/NTU separation]
\label{prop:ntusep}
There are instances whose transferable-utility core is empty while the
$\theta=0$ no-numeraire core is nonempty.
\end{proposition}

\begin{proof}
Take the three-cycle of Theorem~\ref{prop:walras}(b): floors $(1,1,1)$,
pair bids on $\{1,2\},\{2,3\},\{1,3\}$, each with componentwise vector
$(1+\delta/2,\,1+\delta/2)$, $\delta>0$. \emph{TU core empty:} the core
constraints include $x_j+x_{j'}\ge2+\delta$ for each pair; summing the
three gives $2(x_1+x_2+x_3)\ge3(2+\delta)$, i.e.\
$\sum x_j\ge3+\tfrac{3\delta}2>3+\delta=\OPT$, contradicting
$\sum x_j=\OPT$. \emph{$\theta=0$ NTU core nonempty:} allocate the batch
$\{1,2\}$ delivering $(1+\delta/2,\,1+\delta/2)$ and order $3$
individually at $1$. The batch $\{1,3\}$ cannot strictly improve order
$1$ (its componentwise cap is $1+\delta/2$, already attained); likewise
$\{2,3\}$ for order $2$; singletons offer at most the floors. No bid
blocks.
\end{proof}

\begin{remark}[TU-blocked yet NTU-unblocked]\label{rem:tublocked}
The mechanism behind the separation appears already with two orders:
floors $(1,1)$ and a batch of vector $(2g,\,0)$. The batch creates total
surplus $2g>2$, so the floor allocation is TU-\emph{blocked} (a transfer
could compensate order $2$); yet at $\theta=0$ it is NTU-unblocked, since
the batch cannot deliver order $2$ anything componentwise. (The TU core
of this two-order instance is nonempty---e.g.\ $x_1=x_2=g$---which is why
the genuine emptiness separation needs the overlap geometry of the
three-cycle: TU core emptiness, like the LP gap and Walrasian failure, is
a phenomenon of \emph{overlapping} packages.)
\end{remark}

\begin{remark}[A pattern, and an obstruction]\label{rem:removesdev}
This is the third appearance of ``no-numeraire removes deviations'':
VCG's payments are infeasible (Proposition~\ref{prop:vcg}); batch
deviations are dominated at $g=1$ (Lemma~\ref{lem:domination}); and
componentwise blocking is weak enough that the core can be nonempty where
the transferable core is empty (Proposition~\ref{prop:ntusep}). Whether
the $\theta=0$ core is nonempty for \emph{every} instance is open; the
natural route, Scarf's theorem, requires balancedness, and fractional
combinations of indivisible bids are not physically available---the same
overlap geometry that creates the LP integrality gap
(Remark~\ref{rem:botransfers}). We verified nonemptiness by hand on the
three-cycle and conjecture it in general at $\theta=0$.
\end{remark}

% =====================================================================
\section{Winner Determination: Complexity}\label{sec:complexity}
% =====================================================================

\begin{proposition}[The filter is easy]\label{prop:filter}
Given all bids, computing the floors $c_j$ and filtering unfair batched bids
takes time linear in the total bid size.
\end{proposition}

\begin{theorem}[Hardness]\label{thm:nphard}
For every fixed $g>1$, computing $\OPT$ (equivalently $\FAIR$) is NP-hard,
even on instances where the UDP constraint is vacuous and every bid passes
the fairness filter. Membership in NP is immediate.
\end{theorem}

\begin{proof}
Reduction from \textsc{3-Dimensional Matching} \cite{karp1972}: given
disjoint $X,Y,Z$ with $|X|=|Y|=|Z|=q$ and triples $T\subseteq X\times
Y\times Z$, decide whether a perfect matching exists. Create one order per
element of $X\cup Y\cup Z$, each on a distinct directed token pair (UDP
vacuous). Fix $\delta\in(0,3(g-1)]$, possible for any $g>1$. Each element
has an individual bid of value $1$ (so $c_j=1$); each triple $t\in T$ is a
batched bid on its three orders with production vector
$(1+\delta/3,1+\delta/3,1+\delta/3)$, of value $3+\delta\le 3g$
(complementarity factor $\le g$, taking the batching solver's standalone
values to be $1$), delivering $>1$ per order (passes the filter). Any
allocation decomposes into disjoint triples plus individual bids, so
$\OPT=3q+\delta\cdot(\text{max \# disjoint triples in }T)$; hence
$\OPT\ge 3q+\delta q$ iff a perfect matching exists.
\end{proof}

\begin{proposition}[UDP parametrization; FPT]\label{prop:fpt}
Under UDP, all orders on a directed token pair must clear in one winning
bid, so each directed pair is an atomic item. Winner determination has
effective item count $P=$ number of distinct directed pairs, not $n$: it is
solvable in time $2^{O(P)}\cdot\mathrm{poly}(|\mathcal B|)$ ($\mathcal B$ the bid set) by enumeration over pair
subsets (or as weighted set packing over $P$ items), hence polynomial
whenever $P=O(\log n)$, and fixed-parameter tractable in $P$.
\end{proposition}

\begin{remark}[Complexity dichotomy; further standard facts]
\label{rem:complexitydicho}
(i) If every bid covers at most $d$ orders, winner determination is
weighted $d$-set packing; greedy-by-value is a $d$-approximation, and
constant $d$ (deployment: $d\le5$ \todo{verify current CIP-67 batch cap})
admits exact solution in practice. (ii) The no-numeraire constraint creates
no novel equilibrium-computation hardness: the induced exchange economy has
linear single-asset demands (polynomial-time equilibrium); superadditive
production breaks equilibrium \emph{existence}
(Proposition~\ref{prop:walras}) rather than creating a total PPAD-hard
problem; and Nash hardness of the bidding game is generic, not
no-numeraire-specific. Together with Theorem~\ref{thm:nphard} and
Proposition~\ref{prop:fpt}: \emph{the no-numeraire constraint reshapes the
economics (Theorems~\ref{thm:pof}, \ref{thm:incompat}, \ref{thm:posp}), not
the computation.}
\end{remark}
% =====================================================================
\section{Incentives: the Trilemma and the Price of Strategyproofness}
\label{sec:incentives}
% =====================================================================

\subsection{Two impossibilities}

\begin{proposition}[VCG is infeasible {\cite{canidio2024}}]\label{prop:vcg}
No VCG mechanism is feasible under the no-numeraire constraint: the Clarke
payment of a winner may be owed in value the winner cannot produce (an
asset outside its production set), and there is no money in which to pay it.
\end{proposition}

\begin{theorem}[Best-bid fairness is incompatible with strategyproofness]
\label{thm:incompat}
Call a mechanism \emph{non-wasteful} if, whenever some solver's reported
standalone value for an order is positive, the order is allocated; and
\emph{value-respecting} if the order is allocated to a solver whose
reported value for it is maximal. No feasible, non-wasteful,
value-respecting, dominant-strategy incentive-compatible mechanism
guarantees fairness at the best-individual-bid benchmark
$c_j=\max_i(\text{individual bid for }j)$. This holds already for $n=1$.
Under any DSIC mechanism in this class the effective enforceable benchmark
collapses to the runner-up (second-price) level. (Some such richness
condition is necessary: the never-allocating mechanism is vacuously fair
and DSIC.)
\end{theorem}

\begin{proof}
Let $n=1$. Solver $i$ has private deliverable value $v_i$ for the single
order and reports $r_i\le v_i$ (bids are binding; over-reporting is
excluded by settlement, under-reporting is free). By non-wastefulness the
order is allocated; by value-respecting the winner $w$ has maximal report,
and it must deliver some $d\le v_w$ (feasibility). Best-bid fairness
requires $d\ge\max_i r_i$---note the benchmark is \emph{endogenous} to the
reports. Suppose the top producer (value $v^{(1)}$, runner-up value
$v^{(2)}<v^{(1)}$) reports truthfully: it sets the benchmark to $v^{(1)}$,
must deliver $d=v^{(1)}$, and earns $v^{(1)}-d=0$. Deviating to
$r=v^{(2)}+\varepsilon$ it still wins (its report remains maximal), the
benchmark drops to $v^{(2)}+\varepsilon$, and it earns
$v^{(1)}-v^{(2)}-\varepsilon>0$. Hence
truthful reporting is not dominant for any mechanism enforcing the best-bid
benchmark; and in any DSIC mechanism in this class, the enforceable
delivery is bounded by what the runner-up disciplines, i.e.\ the
second-price level.
\end{proof}

\begin{corollary}[Endogenous-floor manipulation at the batch level]
\label{cor:batchmanip}
For $n\ge2$, a batch winner profitably under-reports its \emph{individual}
bids to depress the fairness floors its batch must clear, then
under-delivers; Theorem~\ref{thm:incompat} applied per order formalizes the
two manipulation channels flagged in the CIP-67 documentation.
\todo{cite the specific CoW governance forum / docs passages.}
\end{corollary}

\begin{remark}[The trilemma]\label{rem:trilemma}
Feasibility, best-bid fairness, strategyproofness, and efficiency are not
jointly achievable, and the three natural mechanisms each drop exactly one:
individual second-price auctions are feasible, SP, and (individually)
efficient, but only runner-up-fair; CIP-67 welfare-max-over-fair-bids is
feasible, fair, and efficient, but not SP
(Theorem~\ref{thm:incompat}/Corollary~\ref{cor:batchmanip}); VCG is fair,
SP, and efficient, but infeasible (Proposition~\ref{prop:vcg}). This
organizes the first-price-vs-second-price tradeoff observed in the
$2\times2$ equilibrium analysis of \cite{canidio2024}.
\end{remark}

\subsection{The price of strategyproofness is \texorpdfstring{$\Theta(\log g)$}{Theta(log g)}}
\label{sec:posp}

How much welfare must strategyproofness cost? We answer exactly (up to a
constant) for the canonical single-batch screening core of the problem, to
which we reduce; the composition to general instances is discussed at the
end of the section.

\paragraph{The reduction instance $\mathcal{R}_g$.}
Two orders $\{1,2\}$ on distinct directed pairs. A \emph{public} rival
solver can deliver exactly $1/2$ to each order individually; hence the
floors are $c_1=c_2=1/2$, $C=1$, and the fallback (individual) allocation
has welfare exactly $1$. A distinguished solver (``the bidder'') has
individual capabilities $1/2$ per order and a batched capability producing
$(V/2,\,V/2)$, where $V\in[1,g]$ is its private type; the complementarity
bound $V\le g\cdot(\tfrac12+\tfrac12)$ is exactly the constraint $V\le g$.
The full-information fair optimum is $\FAIR(V)=V$: the batch wins and
delivers everything, which is fair since $V/2\ge1/2$.

A (direct, randomized) mechanism maps the bidder's report to an outcome:
with some probability the batch wins with a delivery vector $(d_1,d_2)$,
$d_j\le V/2$ (feasibility), $d_j\ge 1/2$ (fairness); otherwise the fallback
allocation is implemented (welfare exactly $1$: any allocated order must
receive $\ge1/2$ and no individual capability exceeds $1/2$). Write
$D:=d_1+d_2\in[1,V]$ for the batch delivery, and for report $V$:
\[
p(V):=\Pr[\text{batch wins}],\qquad
x(V):=\E\big[D\cdot\mathbf{1}\{\text{batch wins}\}\big].
\]
The bidder's utility is its retained surplus,
$u(V)=p(V)\,V-x(V)$, and the delivered welfare is
\[
W(V)\;=\;x(V)\;+\;\big(1-p(V)\big).
\]
The mechanism is \emph{IR} if the bidder may abstain (guaranteeing utility
$0$), and DSIC if truthful reporting is dominant. The \emph{price of
strategyproofness} of a mechanism is
$\sup_{V\in[1,g]}\FAIR(V)/W(V)=\sup_V V/W(V)$.

\begin{lemma}[One-sided envelope inequality]\label{lem:envelope}
For any IR, downward-DSIC mechanism on $\mathcal R_g$ (misreports
restricted to $r\le V$ by bindingness) with measurable allocation
probability $p$,
\[
u(V)\;\ge\;u(1)+\int_1^V p(t)\,dt,\qquad u(1)\ge0 .
\]
Consequently
$x(V)=p(V)\,V-u(V)\le p(V)\,V-\int_1^V p(t)\,dt$.
No monotonicity of $p$ is needed (and none is implied: one-sided IC does
not yield the two-sided Myerson machinery).
\end{lemma}

\begin{proof}
Downward IC says that for every $1\le r\le V$,
\[
u(V)\;=\;p(V)V-x(V)\;\ge\;p(r)V-x(r)\;=\;u(r)+p(r)(V-r).
\]
Since $p\ge0$, $u$ is nondecreasing, hence differentiable almost
everywhere (Lebesgue's theorem on the differentiability of monotone
functions). At every differentiability point $t$, applying the displayed
inequality with $V=t+h$, $r=t$ and letting $h\downarrow0$ gives
$u'(t)\ge p(t)$. Therefore
\[
\int_1^V p(t)\,dt\;\le\;\int_1^V u'(t)\,dt\;\le\;u(V)-u(1),
\]
the second inequality because a nondecreasing function dominates the
integral of its a.e.\ derivative. IR at the bottom type gives $u(1)\ge0$,
and the bound on $x$ follows from $x(V)=p(V)V-u(V)$.
\end{proof}

\begin{theorem}[Price of strategyproofness]\label{thm:posp}
On the class $\{\mathcal R_g\}$:
\begin{enumerate}[label=(\roman*),leftmargin=2em]
\item \textbf{(Randomized lower bound.)} Every feasible, fair, IR, DSIC
mechanism (deterministic or randomized) has
\[
\sup_{V\in[1,g]}\frac{V}{W(V)}\;\ge\;\frac{g\ln g}{2(g-1)}\;\ge\;
\frac{\ln g}{2}.
\]
\item \textbf{(Randomized upper bound.)} The uniform geometric
posted-delivery mechanism---draw $k$ uniformly from
$\{0,1,\dots,K\}$ with $K=\lceil\ln g\rceil$, post total delivery
$D=e^k$ split $(D/2,D/2)$, and let the bidder accept iff it can---is
feasible, fair, IR, DSIC, and achieves
$\sup_V V/W(V)\le e\,(\lceil\ln g\rceil+1)$.
\item \textbf{(Deterministic characterization.)} Among deterministic such
mechanisms the optimum is exactly $\sqrt g$, achieved by the posted
delivery $D=\sqrt g$.
\end{enumerate}
Hence on the screening core the price of strategyproofness is
$\Theta(\log g)$, and $\sqrt g$ exactly under determinism. Since
$\{\mathcal R_g\}$ is a subclass of instances, the lower bound (i) holds
verbatim for the full model; the upper bound is lifted to general
instances with single-minded solvers under public benchmarks in
Theorem~\ref{thm:composition}, with the multi-minded case open
(Remark~\ref{conj:composition}).
\end{theorem}

\begin{proof}
\emph{(i)} Suppose $\sup_V V/W(V)\le r$, i.e.\ $W(V)\ge V/r$ for all
$V\in[1,g]$. Since $W(V)=x(V)+1-p(V)$ and $p(V)\le1$,
\[
x(V)\;\ge\;\frac{V}{r}-\big(1-p(V)\big)\;\ge\;\frac{V}{r}-1
\qquad\text{for all }V\in[1,g].
\]
By Lemma~\ref{lem:envelope},
$x(V)\le\psi(V):=p(V)\,V-\int_1^V p(t)\,dt$. Integrate
$\psi(V)/V^2$ over $[1,g]$ and apply Fubini (all integrands are
nonnegative and bounded; $p$ is measurable and bounded by assumption):
\[
\int_1^g\frac{\psi(V)}{V^2}\,dV
=\int_1^g\frac{p(V)}{V}\,dV-\int_1^g p(t)\Big(\frac1t-\frac1g\Big)dt
=\frac1g\int_1^g p(t)\,dt\;\le\;\frac{g-1}{g}.
\]
On the other hand,
\[
\int_1^g\frac{V/r-1}{V^2}\,dV
=\frac{\ln g}{r}-\Big(1-\frac1g\Big).
\]
Combining, $\dfrac{\ln g}{r}\le 2\Big(1-\dfrac1g\Big)$, i.e.\
$r\ge\dfrac{g\ln g}{2(g-1)}\ge\dfrac{\ln g}{2}$.

\emph{(ii)} The mechanism is posted (the menu does not depend on the
report), hence DSIC and IR; accepted deliveries satisfy $D\ge e^0=1$ with
per-order delivery $D/2\ge1/2$ (fair) and are accepted only when
$V\ge D$ (feasible). For type $V\in[1,g]$ pick the integer $k\le K$ with
$e^k\le V<e^{k+1}$ (it exists since $V\le g\le e^K$ when $K=\lceil \ln g
\rceil$). With probability $\frac1{K+1}$ the posted delivery is $e^k\le V$,
which the bidder accepts (retaining $V-e^k\ge0$); all other realizations
contribute welfare $\ge\min(\text{accepted }D,\,1)\ge$ the fallback $1>0$.
Hence
\[
W(V)\;\ge\;\frac{e^k}{K+1}\;\ge\;\frac{V}{e\,(K+1)},
\qquad\text{so}\qquad
\sup_V\frac{V}{W(V)}\;\le\;e\,(\lceil\ln g\rceil+1).
\]
\todo{tighten: acceptance at posted $e^j\le e^k$ for $j<k$ also
contributes; constant improvable, not pursued.}

\emph{(iii)} Let the mechanism be deterministic, so $p(V)\in\{0,1\}$.
First, $p$ is a threshold: if $p(r)=1$ and $p(V)=0$ for some $r<V$, then
type $V$ has truthful utility $0$, while under-reporting $r$ yields utility
$V-x(r)\ge V-r>0$ (feasibility at report $r$ gives $x(r)\le r$),
contradicting downward IC. So there is $T\in[1,g]\cup\{\infty\}$ with
$p=\mathbf 1[V\ge T]$ (up to the boundary point). Second, the accepted
delivery is bounded by $T$: for $T\le r<V$, downward IC gives
$V-x(V)\ge V-x(r)$, i.e.\ $x(V)\le x(r)$; letting $r\downarrow T$ and
using feasibility $x(r)\le r$ yields $x(V)\le T$ for all $V\ge T$, in
particular $x(g)\le T$. Now: if $T>1$, types $V\uparrow T$ have $W(V)=1$,
forcing ratio $\ge T-\epsilon$ for every $\epsilon>0$, and type $g$ has
$W(g)=x(g)\le T$, forcing ratio $\ge g/T$; hence the price is
$\ge\max(T,g/T)\ge\sqrt g$. If $T=1$ (the mechanism always allocates),
$x(g)\le1$ and type $g$ forces ratio $\ge g\ge\sqrt g$. The posted
delivery $D=\sqrt g$ (accept iff $V\ge\sqrt g$) attains
$\sup_V V/W(V)=\max\big(\sup_{V<\sqrt g}V/1,\;
\sup_{V\ge\sqrt g}V/\sqrt g\big)=\sqrt g$.
\end{proof}

\begin{remark}[Reading the bound]\label{rem:pospreading}
Naive individual (second-price) auctions are feasible and SP but lose the
full factor $g$ against the fair optimum. Theorem~\ref{thm:posp} says
randomized posted delivery recovers all but a $\Theta(\log g)$ factor, and
that no DSIC mechanism can do better: \emph{strategyproofness costs
$\log g$, not $g$.} The proof of (i) is intrinsic to the no-numeraire
model: the screened quantity is the solver's deliverable value itself, the
fairness floor enters as a delivery constraint $D\ge C$ rather than a
revenue objective, and welfare---not revenue---is what the kernel
$V^{-2}$ prices. Section~\ref{sec:composition} lifts the upper bound to general
single-minded instances under public benchmarks; the multi-minded
characterization (taxation principle for multi-batch menus) remains the
open core of the program (Remark~\ref{conj:composition}).
\end{remark}

\subsection{Lifting to general instances: composition under public
benchmarks}\label{sec:composition}

Theorem~\ref{thm:posp} characterizes the price of strategyproofness on the
single-batch core. We now lift the upper bound to general instances. The
public-floor model in which we do so is not merely a convenience.
Theorem~\ref{thm:incompat} rules out dominant-strategy truthfulness when
the floor is endogenously set by the same reports that determine the
allocation; Proposition~\ref{prop:online} shows that, under a strong
lifetime-scoped online benchmark, realized future best bids cannot be
implemented by any deterministic online fair mechanism. Together these
results \emph{motivate} the exogenous-benchmark regime: floors fixed
before the strategic allocation problem, for example by a historical or
predetermined reference procedure. (We do not claim deployed mechanisms
are literally of this form---CIP-67's filter is per-auction---only that
the impossibilities above single this regime out as the natural one in
which the strategyproofness question is well-posed.) The composition
theorem below solves the lift in this regime for public single-minded
solvers.

\paragraph{Setting (public $g$-calibrated benchmark, single-minded
solvers).}
Floors $c_j>0$ are public; write $C_S:=\sum_{j\in S}c_j$.
\textbf{(A1)} Every order can always be settled at exactly $c_j$ through a
public reference supply (the benchmark is executable), so the fallback
allocation has welfare $C$ and is fair. \textbf{(A2)} Each solver $i$ is
single-minded: it has one public order set $S_i$ and a private
filter-passing batch, i.e.\ a production vector with $v_{i,j}\ge c_j$
componentwise and private total value $V_i\in[C_{S_i},\,g\,C_{S_i}]$;
singletons are batches with $|S_i|=1$. Each strategic identity submits one
bid, the set $S_i$ is public, and the private type is the \emph{scalar}
$V_i$; reports are binding scalars $r_i\le V_i$, as in
Lemma~\ref{lem:envelope}. Any total $D\in[C_{S_i},V_i]$ is deliverable
with all floors cleared componentwise: if $V_i>C_{S_i}$, set
$\alpha=\frac{D-C_{S_i}}{V_i-C_{S_i}}\in[0,1]$ and deliver
$w_j=c_j+\alpha(v_{i,j}-c_j)$, so that $c_j\le w_j\le v_{i,j}$ and
$\sum_j w_j=D$ (if $V_i=C_{S_i}$, deliver $w_j=c_j=v_{i,j}$). The
mechanism need not observe the vector $v_i$; settlement verifies only
that the delivered vector clears the floors and totals the required
amount. $\FAIR$ denotes the fair-optimal welfare under true values.

\paragraph{The eligibility-posted packing mechanism $M^\ast$.}
(1) Draw $k_i\sim\mathrm{Unif}\{0,\dots,K\}$ independently per solver,
$K=\lceil\ln g\rceil$, and post $P_i:=C_{S_i}\,e^{k_i}$.
(2) Let the eligible set be $E:=\{i:r_i\ge P_i\}$.
(3) Let $A$ be a maximum-weight disjoint packing over $E$ with the
\emph{public} weights $P_i$ (fixed deterministic tie-breaking).
(4) Each $i\in A$ wins $S_i$ and is required to deliver total $P_i$ with
all floors cleared componentwise; by (A2) such a delivery exists whenever
$P_i\le V_i$.
(5) Orders not covered by $A$ settle at floors via the reference supply.

\begin{theorem}[Composition]\label{thm:composition}
Assume: public floors with executable reference supply (A1); public
single-minded solvers with binding scalar reports and $g$-calibrated
filter-passing batches (A2); and tie-breaking among maximum-weight
packings by a fixed deterministic rule depending only on public data and
the realized posts $P$, never on reports beyond eligibility. Then
$M^\ast$ is feasible, fair, IR, and universally truthful
(dominant-strategy truthful for every realization of the posted draws),
and on every instance
\[
\E[W(M^\ast)]\;\ge\;\frac{\FAIR}{e\,(\lceil\ln g\rceil+1)+1}
\;\ge\;\frac{\FAIR}{2e\,(\lceil\ln g\rceil+1)} .
\]
Combined with the embedded lower bound of Theorem~\ref{thm:posp}(i)
(disjoint copies of $\mathcal R_g$ satisfy (A1)--(A2)), the price of
strategyproofness for single-minded solvers under a public $g$-calibrated
benchmark is $\Theta(1+\log g)$, i.e.\ $\Theta(\log g)$ as
$g\to\infty$. $M^\ast$ is an information-theoretic mechanism: the
packing step is NP-hard in general, exactly as winner determination itself
(Theorem~\ref{thm:nphard}).
\end{theorem}

\begin{proof}
\emph{Truthfulness.} Fix the draw realization and $r_{-i}$. The packing
weights $P$ are public and report-independent, and ties are broken by a
fixed deterministic rule depending only on public data and $P$; hence
whether $i\in A$ depends on $r_i$ only through the boolean event
$r_i\ge P_i$. So $i$'s allocation is
$\mathbf 1[r_i\ge P_i]\cdot\chi_i$ with $\chi_i$ (``$i$ is selected when
eligible'') determined by $(k,r_{-i})$ alone, and its required delivery
$P_i$ is also independent of $r_i$. Solver $i$ thus faces a posted
take-it-or-leave-it offer: if $V_i\ge P_i$, truthful reporting secures the
nonnegative option $V_i-P_i$ and under-reporting below $P_i$ only forfeits
it; if $V_i<P_i$, no feasible report wins. Truthfulness is weakly dominant
for every realization. (Over-reporting is excluded by bindingness; even if
permitted, winning with $P_i>V_i$ is infeasible at settlement.)

\emph{Feasibility, fairness, IR.} A winner delivers $P_i\le r_i\le V_i$,
with componentwise floors cleared by the scaling in (A2) since
$P_i\in[C_{S_i},V_i]$; uncovered orders receive exactly $c_j$; the winner's
utility is $V_i-P_i\ge0$.

\emph{Welfare.} Let $A^\ast$ be the fair-optimal allocation, with batch set
$B^\ast$ (values $V_b\in[C_b,gC_b]$ by (A2)) and individual orders settled
at floors, so $\FAIR=\sum_{b\in B^\ast}V_b+C_{\mathrm{ind}}$ where
$C_{\mathrm{ind}}$ is the floor mass of $A^\ast$'s individual orders. For
$b\in B^\ast$ let $\mathcal E_b$ be the \emph{bracket event}
$P_b\in(V_b/e,\,V_b]$: since the grid $\{C_b e^k\}_{k=0}^{K}$ spans
$[C_b,\,C_b e^K]\supseteq[C_b,\,gC_b]$, there is exactly one $k(b)\le K$
with $C_be^{k(b)}\le V_b<C_be^{k(b)+1}$, so
$\Pr[\mathcal E_b]=\tfrac1{K+1}$ (only linearity of expectation is used
below; no independence across $b$ is needed). Pointwise in the draw: the
eligible members of $B^\ast$ form a feasible disjoint packing inside $E$,
so the max-weight packing satisfies
\[
\sum_{i\in A}P_i\;\ge\;\sum_{b\in B^\ast}P_b\,\mathbf 1[b\in E]
\;\ge\;\frac1e\sum_{b\in B^\ast}V_b\,\mathbf 1[\mathcal E_b].
\tag{$\ast$}
\]
Also, every order---covered (componentwise floors) or uncovered
(fallback)---receives at least $c_j$, so $W(M^\ast)\ge C\ge
C_{\mathrm{ind}}$ pointwise. Taking expectations of $(\ast)$ by linearity,
$\E[W]\ge\frac1{e(K+1)}\sum_{B^\ast}V_b$. Writing $a:=e(K+1)$,
$X:=\sum_{B^\ast}V_b$, $Y:=C_{\mathrm{ind}}$, and using
$\max(X/a,\,Y)\ge\frac{X+Y}{a+1}$ (if $X/a\ge Y$ then
$X+Y\le X\frac{a+1}{a}=(a+1)\frac Xa$; otherwise $X+Y\le(a+1)Y$),
\[
\E[W]\;\ge\;\max\!\Big(\frac{X}{a},\,Y\Big)
\;\ge\;\frac{X+Y}{a+1}
\;=\;\frac{\FAIR}{e(\lceil\ln g\rceil+1)+1}. \qedhere
\]
\end{proof}

\paragraph{The multi-minded case: coupling.}
A solver with several private batches is a multi-parameter agent, and
$M^\ast$ breaks: it may withhold a high-weight option to steer the packing
toward a more profitable one. Two results delineate exactly what survives.
The first shows the natural fix---independent posted prices per
option---fails badly: within the posted-menu route, coupling is not a
constant-factor optimization but the structural ingredient that prevents
cherry-picking. The second shows that \emph{coupling} a solver's
prices through a single multiplier restores the logarithmic guarantee
completely, for a single solver of arbitrary multi-minded structure.
(Proposition~\ref{prop:indep} rules out independent geometric posted
menus, not every conceivable uncoupled mechanism.)

\begin{proposition}[Independent pricing fails]\label{prop:indep}
For every $g$ there are single-solver instances satisfying (A1)--(A2)
(with (A2) read per option) on which the mechanism posting independent
geometric prices per option has
$\FAIR/\E[W]=\Omega\big(g/\log^2 g\big)$.
\end{proposition}

\begin{proof}
Sets $S_t=\{0,t\}$ for $t=1,\dots,m$ with a common order $0$;
$c_0=c_t=\tfrac12$, so $C_{S_t}=1$ and any two options conflict. True
values $V_t=g$ for all $t$ (calibrated). With independent draws
$P_t=e^{k_t}$, the solver chooses a minimum posted price among acceptable
options; hence the delivered amount is $D=P_{\min}$ if $P_{\min}\le g$,
and $D=0$ otherwise. In all cases $D\le 1+g\,\mathbf 1\{P_{\min}\ge e\}$.
Since $\Pr[P_{\min}\ge e]=(K/(K+1))^m\le e^{-m/(K+1)}$, choosing
$m=\lceil(K+1)\ln g\rceil$ gives $\Pr[P_{\min}\ge e]\le1/g$, hence
$\E[D]\le 2$, so
$\E[W]\le 2+C=O(\log^2 g)$ while $\FAIR=g+\tfrac{m-1}{2}=\Omega(g)$.
\end{proof}

\begin{theorem}[Coupled menu; single multi-minded solver]
\label{thm:coupled}
Let a single solver hold an arbitrary finite family of filter-passing
options $\{(S,V_S)\}$ satisfying (A1)--(A2) per option, with private
values, and assume execution over disjoint options is \emph{separable}:
the solver can serve any disjoint collection $\mathcal C$ simultaneously,
with value $V(\mathcal C)=\sum_{S\in\mathcal C}V_S$ and each member's
delivery produced as in (A2) (equivalently, each feasible collection may
be regarded as a primitive option). The solver demands any disjoint
collection. The coupled
geometric menu---draw one $k\sim\mathrm{Unif}\{0,\dots,K\}$,
$K=\lceil\ln g\rceil$, set the multiplier $t=e^k$, price each option at
$C_S\,t$, let the solver take any profit-maximizing disjoint collection
(possibly empty) and deliver $C_S\,t$ per taken option componentwise, with
uncovered orders at floors---is universally truthful, feasible, fair, IR,
and on every instance
\[
\E[W]\;\ge\;\frac{\FAIR}{(e-1)\,(\lceil\ln g\rceil+1)+1}.
\]
Hence the price of strategyproofness for a single multi-minded solver
under a public $g$-calibrated benchmark is $\Theta(1+\log g)$, with the
lower bound from the embedded $\mathcal R_g$.
\end{theorem}

\begin{proof}
\emph{Truthfulness, feasibility, fairness, IR.} This is a posted demand
mechanism, not a direct-revelation mechanism: the menu is independent of
any report, and the solver's dominant action is to choose a
utility-maximizing disjoint collection from the posted menu. Every
realization of $t$ is therefore truthful in the taxation/menu sense, and
the randomized mechanism is universally truthful. In any profit-maximizing collection each member
individually has $V_S\ge C_S t$ (dropping a set is feasible and would
otherwise raise profit), so the required delivery is feasible, and
$C_St\ge C_S$ clears all floors componentwise by the scaling of (A2);
uncovered orders settle at floors; the solver's profit is nonnegative.

\emph{Welfare: the profit curve.} Define, for $t\ge1$,
\[
p(t)\;:=\;\max_{\mathcal C\ \mathrm{disjoint}}\ \sum_{S\in\mathcal C}
\big(V_S-C_S\,t\big),
\]
the empty collection included. As a maximum of finitely many affine
functions of $t$, $p$ is convex, piecewise linear, and nonincreasing
(slopes $-C_{\mathcal C}:=-\sum_{S\in\mathcal C}C_S\le0$); and $p(t)=0$
for $t\ge g$ since each term is at most $C_S(g-t)\le0$ by (A2). At any
differentiability point with $p(t)>0$, the solver's delivered total is
$D(t)=t\,C_{\mathcal C^\ast(t)}=t\,(-p'(t))$; at a kink,
$D(t)\ge t\,(-p'_+(t))$ for any tie-breaking, since the right-slope is
the smallest collection weight among maximizers. By convexity $-p'$ is
nonincreasing, so for the geometric grid $t_k=e^k$,
\[
\int_{t_k}^{t_{k+1}}(-p'(t))\,dt\;\le\;(-p'_+(t_k))\,(t_{k+1}-t_k)
\;=\;(e-1)\,t_k\,(-p'_+(t_k))\;\le\;(e-1)\,D(t_k).
\]
Summing over $k=0,\dots,K-1$ and telescoping
($t_K=e^K\ge g$, where $p$ vanishes):
\[
p(1)\;=\;p(t_0)-p(t_K)\;\le\;(e-1)\sum_{k=0}^{K-1}D(t_k),
\qquad\text{so}\qquad
\E[D]\;\ge\;\frac{p(1)}{(e-1)(K+1)} .
\]
\emph{Assembling.} Let $B^\ast$ be the batch collection of the
fair optimum, so $\FAIR=\sum_{B^\ast}V_b+C_{\mathrm{ind}}$. Taking
$\mathcal C=B^\ast$ at $t=1$ gives
$p(1)\ge\sum_{B^\ast}V_b-C_{B^\ast}$, hence
$\FAIR\le p(1)+C_{B^\ast}+C_{\mathrm{ind}}\le p(1)+C$. Pointwise
$W\ge C$ (covered orders clear floors componentwise; uncovered settle at
floors). With $a=(e-1)(K+1)$, $X=p(1)$, $Y=C$ and
$\max(X/a,Y)\ge\frac{X+Y}{a+1}$,
\[
\E[W]\;\ge\;\max\Big(\frac{p(1)}{a},\,C\Big)
\;\ge\;\frac{p(1)+C}{a+1}\;\ge\;\frac{\FAIR}{(e-1)(K+1)+1}.\qedhere
\]
\end{proof}

\begin{remark}[Why coupling works]\label{rem:coupling}
A single multiplier makes the solver's comparison \emph{across its own
options price-neutral}: all margins shrink together as $t$ grows, so the
demanded collection traces the slope of the convex profit curve, and the
geometric grid converts the telescoped curve into delivered welfare. With
independent prices the solver cherry-picks its cheapest own option
(Proposition~\ref{prop:indep}). The coupled constant $(e-1)$ even
improves on the $e$ of Theorem~\ref{thm:composition}. Note also that the
coupled menu requires \emph{no reports at all}---it is a pure demand
mechanism, the taxation-principle object made literal.
\end{remark}

\begin{corollary}[Territorial multi-agent instances]\label{cor:territorial}
If the option families of distinct solvers touch pairwise-disjoint order
blocks---i.e.\ $\big(\bigcup_{S\in\mathcal F_i}S\big)\cap
\big(\bigcup_{S\in\mathcal F_{i'}}S\big)=\varnothing$ for $i\ne
i'$---running the coupled menu independently per solver gives a
universally truthful, fair, feasible mechanism with
$\E[W]\ge\FAIR/\big((e-1)(\lceil\ln g\rceil+1)+1\big)$; the price of
strategyproofness on this class is $\Theta(1+\log g)$.
\end{corollary}

\paragraph{Multi-agent multi-minded: a positive result and a class
barrier.}
The remaining cell is multi-agent \emph{and} multi-minded. We give a
truthful mechanism whose guarantee degrades linearly in the number $N$ of
strategic solvers---which in deployment is a small constant---and then
prove that within the natural class of serial demand mechanisms this
$N$-dependence is unavoidable, already for single-minded agents.

\begin{theorem}[Random-priority coupled menus]\label{thm:randprio}
Under (A1)--(A2) with separable execution, the mechanism $M_{\mathrm{rp}}$
---draw a uniformly random service order over the $N$ solvers and
independent multipliers $t_i=e^{k_i}$, $k_i\sim\mathrm{Unif}\{0,\dots,K\}$,
$K=\lceil\ln g\rceil$; serve solvers in order, each facing the posted
coupled menu of its options fully contained in the remaining orders at
prices $C_S\,t_i$ and taking any profit-maximizing disjoint collection,
delivering $C_S\,t_i$ per taken option componentwise; uncovered orders at
floors---is feasible, fair, IR, and universally truthful, and on every
instance
\[
\E[W]\;\ge\;\frac{\FAIR}{N\,(e-1)(\lceil\ln g\rceil+1)+1}.
\]
(Each solver is served once---no multi-identity/sybil strategies---and has
no externalities over successors' outcomes; the demand step is
information-theoretic, as in Theorem~\ref{thm:coupled}.) In particular,
when the number of strategic solvers is treated as a deployment-scale
constant, the price of strategyproofness in the multi-agent multi-minded
model is $\Theta_N(1+\log g)$.
\end{theorem}

\begin{proof}
\emph{Truthfulness.} Solver $i$'s menu is the pair (remaining order set,
its own posted prices). The remaining set is determined by its
predecessors' demands, hence by their values, never by $i$'s behavior; and
$i$'s demand affects only its successors, whose outcomes do not enter
$i$'s utility. So every realization presents each solver a posted demand
menu, and the mechanism is universally truthful; feasibility, fairness,
and IR are as in Theorem~\ref{thm:coupled}.

\emph{Welfare.} Condition on solver $i$ being served first (probability
$1/N$): the full order set is available and $t_i$ is independent, so the
profit-curve telescope of Theorem~\ref{thm:coupled} applies verbatim:
$\E[\mathrm{del}_i\mid i\text{ first}]\ge p_i(1)/((e-1)(K+1))$, with
$p_i$ computed on the full set. Other positions contribute nonnegatively,
so $\E[\sum_i\mathrm{del}_i]\ge\frac1N\sum_i p_i(1)/((e-1)(K+1))$. Since
each solver's $B^\ast$-batches form a disjoint collection on the full set,
$\sum_i p_i(1)\ge\sum_{b\in B^\ast}V_b-C_{B^\ast}$. Pointwise $W\ge C$.
The $\max(X/a',Y)\ge\frac{X+Y}{a'+1}$ combination with
$a'=N(e-1)(K+1)$ gives the bound.
\end{proof}

\begin{theorem}[Serial demand barrier]\label{thm:serialbarrier}
Call a mechanism a \emph{serial demand mechanism} if it serves identities
in a report-independent (possibly random) identity order $\sigma$---which
may depend on public data, including the identities' public option
geometries, but is drawn \emph{ex ante}: before, and independently of,
all multiplier draws and all realized demands (no adaptive
reordering)---drawing for each identity a multiplier $t_i$ from an
arbitrary report-independent distribution on $[1,\infty)$
(identity-dependent allowed; the $t_i$ are independent across identities
and independent of $\sigma$), and lets each agent at its turn take
a profit-maximizing disjoint collection of its available options at
coupled posted prices $C_S\,t_i$. (The class contains
$M_{\mathrm{rp}}$.) Assume $g\ge2$ (smaller $g$ is absorbed into the
constants). Every serial demand mechanism has worst-case ratio
$\FAIR/\E[W]=\Omega\big(\min(N,g)\big)$, on instances in which every
agent holds two options.
\end{theorem}

\begin{proof}
\emph{Avoiding atoms.} The $N$ multiplier distributions have countably
many atoms in total, so the adversary may fix
$\rho\in\big(\max\{1,g/2\},\,g\big]$ that is an atom of none of them
(the interval is uncountable for $g\ge2$); all values below have ratio at
most $\rho\le g$ (calibrated), and $\rho>1$ strictly.

\emph{The symmetrized swarm.} $N$ identities. Orders $o_1,\dots,o_N$ with
floors $\varepsilon:=\frac{L}{N\rho}$ and $u_1,\dots,u_L$ with floors
$1$; total floor mass $C=N\varepsilon+L\le2L$. \emph{Every} identity $i$
has the \emph{same} public option family
$\{\{o_i\},\,G\}$, where $G$ is the set of all orders, $C_G=C\le2L$.
The adversary places the \emph{batcher} at an identity $\beta$ chosen
after seeing the mechanism; private values and production vectors are:
batcher, $V_\beta(G)=\rho\,C_G$ with vector $(\rho c_j)_{j\in G}$ and
$V_\beta(\{o_\beta\})=\varepsilon$ with vector $(\varepsilon)$; every
other identity (\emph{blocker}) $s$, $V_s(G)=C_G$ with vector
$(c_j)_{j\in G}$ and $V_s(\{o_s\})=\rho\varepsilon$ with vector
$(\rho\varepsilon)$. All options are filter-passing and
$\rho$-calibrated. Crucially, the \emph{public} geometry is identical
across identities and across placements of the batcher; only private
values move, so any report-independent order distribution---geometry-aware
or not---is fixed independently of where the batcher sits.

\emph{Demand behavior.} Fix $t\ne\rho$ (almost surely). A blocker $s$
with $t<\rho$ has singleton profit $\varepsilon(\rho-t)>0$ (at $t=1$ this
is $\varepsilon(\rho-1)>0$ since $\rho>1$) and $G$-profit
$C_G(1-t)\le0$; since $o_s\in G$, the zero-or-negative $G$ conflicts with
the strictly profitable singleton, so the profit-maximizing collection is
exactly $\{o_s\}$---taking $o_s$ and blocking $G$ for all successors.
With $t>\rho>1$ both options are strictly unprofitable
and the blocker is silent. So a blocker is non-silent iff $t_s<\rho$; set
$q_s:=\Pr[t_s>\rho]$. The batcher with $t<\rho$ strictly prefers $G$
(profit $C_G(\rho-t)>0\ge\varepsilon(1-t)$) and delivers $C_G\,t$; with
$t>\rho$ it is silent. If $G$ is unavailable (blocked by a predecessor),
the batcher can take at most its singleton at nonpositive profit,
contributing at most the floor mass already charged in $C$.

\emph{Accounting.} Blockers' total delivery is at most
$N\varepsilon\rho=L$; pointwise $W\le\mathrm{del}_{\mathrm{batcher}}
+L+C\le\mathrm{del}_{\mathrm{batcher}}+3L$. By independence of
$t_\beta$ from the order and the other draws,
\[
\E[\mathrm{del}_{\mathrm{batcher}}]\;\le\;
C_G\cdot\rho\,(1-q_\beta)\cdot
\Pr[\text{all predecessors of }\beta\text{ silent}].
\]
\emph{Telescoping identity.} For any realization of the order $\sigma$,
\[
\sum_{i}(1-q_i)\prod_{s\ \mathrm{precedes}\ i}q_s
\;=\;1-\prod_{s}q_s\;\le\;1
\]
(the $i$-th term is the probability that $i$ is the first non-silent
identity under independent silences). Taking expectation over $\sigma$ and
summing over identities, some $i^\ast$ satisfies
$(1-q_{i^\ast})\cdot\Pr[\text{predecessors of }i^\ast\text{ silent}]
\le1/N$. Placing the batcher at $\beta=i^\ast$:
$\E[\mathrm{del}_{\mathrm{batcher}}]\le2\rho L/N$, hence
$\E[W]\le3L+2\rho L/N$, while $\FAIR\ge V_\beta(G)=\rho C_G\ge\rho L$,
so
\[
\frac{\FAIR}{\E[W]}\;\ge\;\frac{\rho}{3+2\rho/N}
\;=\;\Omega\big(\min(N,\rho)\big)\;=\;\Omega\big(\min(N,g)\big),
\]
using $\rho\ge g/2$.
\end{proof}

\begin{corollary}[Option-geometry-blind subclass; single-minded agents]
\label{cor:geoblind}
If the service-order distribution is additionally
\emph{option-geometry-blind}---it may depend on $N$, $g$, the floor
system, and fixed identity priorities, but not on the identities' option
geometries or covered-set sizes---the $\Omega(\min(N,g))$ bound holds on
instances where every agent is \emph{single-minded}: drop the ratio-one
decoy options from the symmetrized swarm (batcher keeps only $G$, blockers
only their singletons); the public geometries then differ across
identities, but a geometry-blind order cannot exploit this, and the
same telescoping argument applies verbatim. In summary: the wide-class
lower bound uses two-option agents with identical public geometry; for
option-geometry-blind serial mechanisms, the same bound already holds on
single-minded agents.
\end{corollary}

\begin{remark}[The cell, pinned by three theorems]\label{rem:pinned}
Within the option-geometry-blind subclass the barrier already bites for
\emph{single-minded} agents (Corollary~\ref{cor:geoblind}), which shows
the packing structure of Theorem~\ref{thm:composition} is essential
relative to that subclass---packing is not a serial demand mechanism. For
multi-minded agents the situation is now pinned from three sides: packing
is non-truthful (the withholding manipulation of
Section~\ref{sec:composition}); serial demand---even geometry-aware---is
truthful but pays $\Omega(\min(N,g))$ from floor-blocking
(Theorem~\ref{thm:serialbarrier}); and random-priority coupled menus
achieve $O(N\log g)$ (Theorem~\ref{thm:randprio}), matching the barrier
within the class up to a logarithm for $N\lesssim g/\log g$. The sharpened
open question: does \emph{any} universally truthful, fair, feasible
mechanism achieve an $N$-independent $\mathrm{polylog}(g)$ guarantee for
multi-agent multi-minded solvers? By Theorem~\ref{thm:serialbarrier},
such a mechanism must either use a global packing structure or otherwise
break the serial demand format. We also recall that floors here are the
public exogenous benchmark: a blocker's singleton value does not raise any
floor, in contrast with the endogenous-floor model of
Theorem~\ref{thm:incompat}.
\end{remark}

\begin{remark}[Benchmark calibration]\label{conj:composition}
The parameter entering all results of this subsection is the
\emph{effective spread}
relative to the public benchmark, $G_{\mathrm{eff}}:=\max V_S/C_S$, which
equals the complementarity factor $g$ only when the benchmark tracks
standalone capability; when historical floors drift downward,
$G_{\mathrm{eff}}$---and hence the $\log$ factor---grows, which is
precisely the drift phenomenon the rate-limited benchmark of
Section~\ref{sec:online} manages. \todo{Phase-2: quantify calibration drift.}
\end{remark}

% =====================================================================
\section{Calibrating \texorpdfstring{$g$ and $P$}{g and P} on Deployed Auctions}\label{sec:empirics}
% =====================================================================

Every bound in this paper is parameterized by two observables: the
complementarity factor $g$ (Definition~\ref{def:g}) and the directed-pair
count $P$ (Proposition~\ref{prop:fpt}). Both are estimable from public CoW
Protocol auction data, making the theory directly falsifiable on the
deployed mechanism. \todo{This section is a scaffold; populate from the
companion dataset.}

\paragraph{Estimating $g$.}
For each settled auction and each submitted batched bid $(i,S,y)$, the
ratio
\[
\widehat\gamma_{i,S}\;:=\;\frac{V_i(S)}{\sum_{j\in S}V_i(\{j\})}
\]
is observable whenever solver $i$ also submitted standalone bids for each
$j\in S$ (and is a lower bound on the instance's $g$ in general, since
solvers need not reveal full standalone capability). \todo{(a) report the
empirical distribution of $\widehat\gamma$ across auctions/solvers; (b)
report coverage: fraction of batched bids for which all standalone
components are observed; (c) discuss the resulting censoring direction
(estimates are conservative).}

\paragraph{Estimating $P$.}
$P$ is directly observable per auction batch. \todo{(a) distribution of $P$
per auction; (b) empirical runtime of exact winner determination vs $P$,
verifying the $2^P$ envelope of Proposition~\ref{prop:fpt}; (c) fraction of
auctions with $P\le 20$ (say), where exact solution is trivially in reach.}

\paragraph{Where deployment sits on the $\PoF$ surface.}
Combining $\widehat g$ with conversion-cost estimates $\widehat\theta$ from
on-chain AMM depth \todo{methodology from companion paper}, place deployed
markets on the $\min(g,1/\theta)$ surface of Theorem~\ref{thm:pof}: the
theory predicts which markets are in the complementarity-limited regime
($\theta<1/g$) versus the liquidity-limited regime ($\theta>1/g$), with
directly observable consequences for how much welfare the fairness filter
forgoes. \todo{one figure: empirical $(\widehat\theta,\widehat g)$ scatter
over the $\PoF$ heatmap.}

% =====================================================================
\section{Discussion: the Online, Repeated Setting}\label{sec:online}
% =====================================================================

Trade-intent auctions recur every few seconds, and the fairness benchmark
is, in deployment, partly \emph{historical}. We record what worst-case
analysis says about why, and what it cannot yet say.

\paragraph{The online model.}
Some care is needed in defining online fairness: if floors were computed
over an unbounded future stream, \emph{no} mechanism---batched or
individual---could ever fairly settle anything, since any finite delivery
can be undercut later. To isolate the online obstruction, we study a
\emph{strong lifetime-scoped benchmark}: the floor of an order is the best
standalone opportunity that would arrive during its validity window. We
stress this is not the per-auction CIP-67 filter itself; rather, it
formalizes the stronger ex-post guarantee an online mechanism would need
in order to be fair against later quotes arriving before expiration. The
impossibility below then explains why any deployable online notion of
fairness must use a predetermined or historical benchmark instead of the
realized future best bid.

Formally: each order $j$ is alive on an interval $[a_j,d_j]$; bids
naming $j$ arrive adversarially during $j$'s lifetime and remain available
while all their orders are alive; the floor $c_j$ is the best standalone
bid for $j$ over its \emph{entire lifetime}. Bid-arrival streams are
interpreted as potential opportunities for the order during its validity
window and are fixed independently of the mechanism's settlement times:
settling early does not erase the counterfactual bids used to define the
lifetime benchmark. Settlement is per-bid, atomic,
and irrevocable, and a bid can be settled only while all its orders are
alive. An online mechanism observes arrivals in real time, has no knowledge
of the future, and must be fair on \emph{every} instance: each settled
order $j$ must receive $w_j\ge c_j$ against the realized (full-lifetime)
floor. Note that settling order $j$ individually \emph{at its deadline}
$d_j$ with the best standalone bid seen is always fair---by $d_j$ all of
$j$'s bids have arrived and the settlement delivers exactly $c_j$---so
individual-only allocation is safely implementable online. The content of
the next proposition is that nothing more is.

\begin{proposition}[Online fairness forces individual-only allocation]
\label{prop:online}
In the model above, any deterministic online mechanism that guarantees
best-bid fairness on every instance never settles a batched bid containing
an order whose deadline strictly exceeds the settlement time. Consequently,
on pairwise-distinct-deadline instances its allocation is individual-only
and has welfare at most $C$. There are such instances whose offline fair
optimum is $gC$, while the settle-at-deadline individual mechanism is fair
and guarantees ratio at most $g$ on every instance. Hence the deterministic
competitive ratio against the offline fair optimum is exactly $g$.
\end{proposition}

\begin{proof}
\emph{No exposed batch is ever settled.}
Let $M$ be deterministic and fair on every instance. Suppose, toward a
contradiction, that on some instance prefix $\pi$, $M$ settles at time
$\tau$ a batched bid $(i,S,y)$ delivering $w_{j}$ to each $j\in S$, with
some $j^\ast\in S$ satisfying $d_{j^\ast}>\tau$. Consider the continuation
of $\pi$ in which a fresh solver (with no other bids, so
Definition~\ref{def:g} imposes nothing) submits at time
$\tau'\in(\tau,d_{j^\ast})$ a standalone bid for $j^\ast$ of value
$w_{j^\ast}+1$. This is a legal instance; on it, the realized floor
satisfies $c_{j^\ast}\ge w_{j^\ast}+1>w_{j^\ast}$, and the settlement of
$j^\ast$ is already irrevocable, so $M$ is unfair on this
instance---contradiction. Since a batch can only be settled while all its
orders are alive, i.e.\ at $\tau\le\min_{j\in S}d_j$, on any instance whose
deadlines are pairwise distinct \emph{every} batch with $|S|\ge2$ has an
order with $d_{j^\ast}>\tau$ at any admissible settlement time; hence $M$
settles no batch at all on such instances. Every settled order is then
served by a standalone bid, whose delivery is at most its value, which is
at most $c_j$; so $M$'s welfare is at most $\sum_j c_j=C$.

\emph{The hard instance.}
Fix $q\ge1$ and $\varepsilon>0$. Orders $1,\dots,2q$ on distinct directed
pairs, grouped into pairs $S_k=\{2k-1,2k\}$; all orders arrive at time $0$;
deadlines are pairwise distinct (order $m$ has $d_m=1+m\varepsilon$). At
time $0$ each order receives a standalone bid of value $1$ (so $c_j=1$,
$C=2q$), and each pair $S_k$ receives a batched bid of value $2g$ from a
solver with standalone values $1$ per order (the complementarity bound
$2g\le g(1+1)$ holds with equality), delivering $(g,g)$---which clears
both floors, so the bid passes the filter. No further bids arrive. Because
$M$ is deterministic and must be fair on \emph{all} instances, including
the undercut continuations above which share every prefix of this instance
up to any prospective settlement time, $M$ settles no batch here; its
welfare is at most $C=2q$. The offline fair optimum settles all $q$
batches: welfare $2gq$, each order receiving $g\ge1$. Hence the
competitive ratio is at least $2gq/2q=g$.

\emph{Matching upper bound.}
The settle-at-deadline individual mechanism is fair and achieves exactly
$C$ on every instance, while the offline fair optimum is at most
$\OPT\le gC$ (Theorem~\ref{thm:pof}, upper bound, part 1); so its
competitive ratio is at most $g$. The ratio is therefore exactly $g$.
\end{proof}

\paragraph{Does randomization beat $g$ online?}
The answer splits the online problem into two regimes with different
prices. Against \emph{realized lifetime floors}, randomization does not
help at all; once floors are \emph{static and public}, the residual
online problem is classical one-way search over the spread $g$
\cite{elyaniv2001}, with deterministic price $\Theta(\sqrt g)$ and
randomized price $\Theta(\log g)$. Proposition~\ref{prop:online}'s $g$
is thus the price of floor \emph{uncertainty}, not of search.

\begin{theorem}[Randomization does not help against realized floors]
\label{thm:onlinerand}
In the lifetime-scoped benchmark model with almost-sure fairness---every
realization of the mechanism's randomness must respect every order's
realized lifetime floor---the randomized competitive ratio is exactly
$g$, matching Proposition~\ref{prop:online}.
\end{theorem}

\begin{proof}
The upper bound is the deterministic one. For the lower bound, call a
batch settlement \emph{exposed} if some covered order's deadline lies in
the future, so that a later standalone bid can still raise its realized
floor. If a mechanism settles an exposed batch with positive probability
on some observation prefix, fix that prefix and append a higher
standalone bid on a later-deadline covered order: the mechanism's
behavior on the shared prefix is identical (it is online), so on the
appended instance it has, with positive probability, made an irrevocable
settlement below a realized floor---not almost-surely fair. Hence an
almost-surely fair mechanism settles exposed batches with probability
$0$, and on the staggered-deadline hard family of
Proposition~\ref{prop:online}, where every available batch is exposed,
it is individual-only almost surely: expected welfare at most $C$
against $\FAIR=gC$.
\end{proof}

\begin{theorem}[Static public floors: online search over the spread]
\label{thm:onlineub}
Call an instance \emph{static single-contention} if all orders are
present from time $0$ with public floors (total mass $C$ known and
fixed---no future standalone bid moves them), offers arrive online and
are fleeting, and offers pairwise conflict. Define the \emph{stopped
welfare} of a fair offer $(S,v)$ as $q:=V+C-C_S\in[C,gC]$ ($V\ge C_S$
gives the lower end; $V\le gC_S$ and $C_S\le C$ the upper). Then:
(a) the deterministic competitive ratio is $\Theta(\sqrt g)$---accept
the first offer with $q\ge\sqrt g\,C$ for the upper bound; for the lower
bound the adversary shows $q=\sqrt g\,C$ and continues with $gC$ iff it
is accepted; (b) the mechanism that draws
$k\sim\mathrm{Unif}\{0,\dots,K\}$, $K=\lceil\ln g\rceil$, and accepts
the first fair offer with $q\ge e^kC$ achieves
\[
\E[W]\;\ge\;\frac{\FAIR}{e\,(\lceil\ln g\rceil+1)},
\]
and combined with Theorem~\ref{thm:onlinelb} below, the randomized
competitive ratio of the class is $\Theta(\log g)$.
\end{theorem}

\begin{proof}
(a) If the maximal stopped welfare $q^\ast\ge\sqrt g\,C$, the threshold
mechanism accepts some offer with $q\ge\sqrt g\,C$, ratio
$q^\ast/q\le gC/(\sqrt g\,C)=\sqrt g$; otherwise it accepts nothing and
$W=C$, ratio $q^\ast/C<\sqrt g$. The matching lower bound is the
standard two-step adversary. (b) $\FAIR=q^\ast$ (the fallback is already
inside $q$; if no offer is useful, $q^\ast=C$ and $W=C$ matches). With
probability $1/(K+1)$ the draw satisfies $e^kC\le q^\ast<e^{k+1}C$;
under that draw the mechanism accepts some offer---at latest the
offline's---with $q\ge e^kC>q^\ast/e$. Hence
$\E[W]\ge q^\ast/(e(K+1))$. The known anchor $C$---the fallback floor
mass---is what compresses the search range to ratio $g$; with
dynamically arriving orders $C$ grows online and the reduction breaks,
which is why the class is scoped as stated.
\end{proof}

\begin{theorem}[Randomized lower bound for static search]
\label{thm:onlinelb}
Every randomized online fair mechanism, already on static
single-contention instances, has competitive ratio $\Omega(\log g)$
against an oblivious adversary (Yao's principle).
\end{theorem}

\begin{proof}
Order $o$ has public floor $c$ anchored by a standing reference
executable at the deadline; fleeting batch offers
$(\{o,\mathrm{aux}_i\},(x_i,\varepsilon))$, $x_i=e^ic$,
$i=1,\dots,K'$, $K'=\lfloor\ln g\rfloor$ (so $x_{K'}\le gc$ and
Definition~\ref{def:g} holds with the batching solver's standalone value
$c$ on $o$), arrive in increasing order, each available only at its
arrival instant; batches do not move the standalone benchmark, so floors
are static. All offers share $o$: accepting kills the future, rejecting
all yields the fallback $c$. On this family a deterministic mechanism is
exactly a stopping threshold $\tau$ (the only observable is the count of
offers so far; any plan after a first acceptance is mooted, and early
individual settlement of $o$ is the variant $\tau=\infty$ with a smaller
fallback). Hard distribution: show the first $M$ offers with
$\Pr[M\ge i]=e^{-i}$, $i\le K'$, terminal atom $\Pr[M=K']=e^{-K'}$.
Then for every $\tau\le K'$,
$\E[\mathrm{ALG}_\tau]=e^{-\tau}e^{\tau}c+(1-e^{-\tau})c
+O(K'\varepsilon)\le2c+o(c)$, and $\E[\mathrm{ALG}_\infty]=c$, while
$\E[\FAIR]=c\,\E[e^M]+O(c)=\Omega(cK')$. By Yao's principle every
randomized mechanism has ratio $\Omega(K')=\Omega(\log g)$.
\end{proof}

\begin{proposition}[Ratio-threshold greedy pays $\Omega(g)$ online]
\label{prop:onlinegreedy}
Any mechanism that draws a sequence-independent random threshold $T$ and
greedily accepts every arriving fair offer with value-to-floor ratio
$V/C_S\ge T$ has competitive ratio $\Omega(g)$, already with static
public floors: at time $1$ a blocker $(\{o\},g\varepsilon)$ of ratio $g$
arrives (floor of $o$ is $\varepsilon$), at time $2$ the batch
$(\{o\}\cup\{u_1,\dots,u_L\},\,g(\varepsilon+L))$, componentwise fair,
of ratio at most $g$. Every realization $T\le g$ accepts the blocker,
irrevocably settling $o$ and killing the batch (welfare at most
$L+g\varepsilon$); every $T>g$ accepts nothing (welfare $\approx L$);
with $\varepsilon\ll L/g$, $\FAIR\approx gL$.
\end{proposition}

\begin{remark}[Two online prices; kill potential]
\label{rem:killpotential}
The split is the online strengthening of the paper's causal chain:
realized lifetime floors are too strong a benchmark---even randomization
cannot safely batch against them (Theorem~\ref{thm:onlinerand})---which
is one more reason the deployed design fixes \emph{public, historical}
floors; once it does, the residual online problem is one-way search over
the spread, with tight randomized price $\Theta(\log g)$ on the
single-contention class. Beyond single contention, the obstruction to a
general $O(\log g)$ mechanism is the \emph{kill-potential asymmetry}:
settling an order at any value risks blocking future batches of
unboundedly larger floor mass attached to it
(Proposition~\ref{prop:onlinegreedy} is the swarm of
Theorem~\ref{thm:serialbarrier} in online clothing), so beating greedy
requires rejecting high-ratio low-surplus offers---plausibly via
surplus-budgeted rejection, which we leave open, conjecturing
$\Theta(\log g)$ for general static-floor instances. The geometric-level
hardness is the same scale-free family that appears offline in the
screening core $R_g$ (Theorem~\ref{thm:posp}): one family governs the
informational and the incentive price of the spread $g$.
\end{remark}

\subsection{The contamination game: pricing benchmark manipulation}
\label{sec:contamination}

The rate-limited historical benchmark invites a final strategic
question: what does it cost to \emph{manipulate} the floors themselves?
We give a deliberately reduced-form model whose parameters are directly
estimable from deployed data (Section~\ref{sec:empirics}).

\paragraph{Model.}
One directed pair; the honest benchmark $b^\ast$ is a lower-tail
$q$-quantile of each round's standalone quotes (honest mass $m$, density
$f$ at the quantile). An attacker executing volume $v$ per round injects
contamination mass $a_t\in[0,\bar a]$ of low quotes at convex cost
$K(a)=\kappa_1a+\kappa_2a^2/2$, $\kappa_1>0$; linearizing the quantile,
the raw contaminated benchmark is $b^\ast-\beta a_t$ with
$\beta=1/(mf(b^\ast))$ (valid for $a\ll qm$). The published benchmark
obeys the downward rate limit $r$ with instant upward recovery:
$b_t=\min\big(b^\ast,\max(b^\ast-\beta a_t,\,(1-r)b_{t-1})\big)$. The
attacker's per-round payoff is $v(b^\ast-b_t)-K(a_t)$: obligation
reduction on executed volume, less injection cost. (Memoryless window
and a mechanical benchmark; the attacker is a single decision-maker.)

\begin{proposition}[Steady-state bias; the participation threshold]
\label{prop:contbias}
Assume $\kappa_2>0$. If $v\beta\le\kappa_1$ the unique best response is
$a\equiv0$: a marginal injection cost at least $v\beta$ deters
contamination completely (at the boundary $v\beta=\kappa_1$, any $a>0$
loses $\kappa_2a^2/2$). Otherwise the optimal constant attack is
$a^\ast=\min\big(\bar a,(v\beta-\kappa_1)/\kappa_2\big)$, with steady
relative bias $\mu:=\beta a^\ast/b^\ast$ and per-round rent
$\pi(a^\ast)$, which equals $(v\beta-\kappa_1)^2/(2\kappa_2)$ when the
cap is slack ($\bar a\ge(v\beta-\kappa_1)/\kappa_2$) and
$(v\beta-\kappa_1)\bar a-\kappa_2\bar a^2/2$ when it binds; of this,
$K(a^\ast)$ is deadweight and $v\beta a^\ast$ a transfer from traders.
(With $\kappa_2=0$ the problem is bang-bang: $a^\ast=\bar a$ iff
$v\beta>\kappa_1$.)
Vulnerability is governed by $\beta=1/(mf(b^\ast))$: thin honest mass or
thin quantile density makes a pair cheap to bias.
\end{proposition}

\begin{proposition}[Constant attack is optimal; the rate limit is a ramp
tax]\label{prop:contramp}
The dynamics give $b_t\ge b^\ast-\beta a_t$ pointwise, so for every
horizon $T$, every strategy earns at most
$\sum_t[v\beta a_t-K(a_t)]\le T\,\pi(a^\ast)$: bursts cannot improve on
the steady-state per-period rent, and a constant attack $a^\ast$ attains
this bound asymptotically, after the rate-limit ramp. The rate limit
does not change the steady state; its exact role is a \emph{ramp tax}:
reaching relative bias $\mu$ from $b^\ast$ requires at least
$\lceil\ln\frac1{1-\mu}/\ln\frac1{1-r}\rceil$ rounds of descent
($T_\mu\approx\mu/r$ for small parameters). For an attacker who
optimizes ramp injections (hugging the rate-limit path with the minimal
injection $a_t=\Delta_t/\beta$, $\Delta_t=\min\{\Delta^\ast,
b^\ast(1-(1-r)^t)\}$), the exact net rent lost per restart is
\[
L_{\mathrm{net}}\;=\;\sum_{t=1}^{T_\mu}
\Big[\pi(a^\ast)-\pi\big(\Delta_t/\beta\big)\Big]
\;\approx\;\frac{\pi(a^\ast)\,\mu}{3r}
\]
in the small-$(r,\mu)$ interior approximation. Since recovery is
instant, any pause is a full restart, so rate limits leave the
persistent bias unchanged but tax episodic attacks through the finite
ramp path; the knob $r$ trades honest-price responsiveness against
contamination rents on this quantified frontier. (The forgone
\emph{gross} benefit per ramp is
$\approx v\Delta^\ast\mu/(2r)$; the net figure above accounts for the
injection savings while ramping.)
\end{proposition}

\begin{remark}[Detection tempers; only fees deter]\label{rem:detection}
Add per-round detection probability $p(a)$, increasing with $p(0)=0$,
detection ending the game; discount $\delta$. The constant-policy value
is $V(a)=\pi(a)/\big(1-\delta(1-p(a))\big)$ with
$\pi(a)=v\beta a-K(a)$. Its sign is the sign of $\pi(a)$, so detection
never changes \emph{participation} under stationary policies; and at
$a^\ast$, $\frac{d}{da}V\propto-\pi(a^\ast)\,\delta p'(a^\ast)<0$, so it
strictly reduces \emph{intensity}. Deterrence is the cost floor's job:
fees, stake, or bonding of magnitude $\kappa_1\ge v\beta$
(Proposition~\ref{prop:contbias}). The scope is a detection technology
that only terminates or resets the continuation value; a fine, slashing
loss, or identity-exclusion cost enters as an additional penalty and can
change participation.
\end{remark}

\begin{remark}[The $G_{\mathrm{eff}}$ weld; estimands]\label{rem:weld}
Steady contamination scales the pair's floors by $1-\mu$. For any
guarantee whose only calibration input is the public-floor spread
$G=\max_iV_i/C_{S_i}$---the posted-menu logarithms of
Section~\ref{sec:incentives} and the online search bounds above---the
replacement is $G_{\mathrm{eff}}=G/(1-\mu)$, making the calibration
remark (Remark~\ref{conj:composition}) quantitative: contamination
spends exactly like spread inflation in every floor-calibrated
guarantee. The replacement does \emph{not} touch structural production
complementarity (the Walrasian frontier of
Theorem~\ref{prop:walras}), directed-pair complexity, or the
produced-welfare price of anarchy; and for the coverage law
(Theorem~\ref{thm:coverage}) the calibration is joint:
$(g,\alpha)\mapsto\big(g/(1-\mu),\,\min\{1,\alpha/(1-\mu)\}\big)$,
since coverage is itself floor-calibrated. This gives
Section~\ref{sec:empirics} its estimands: $\widehat\beta$ (quantile
sensitivity), $\widehat\kappa$ (injection cost), $\widehat\mu$, the
implied $\widehat G_{\mathrm{eff}}=\widehat g/(1-\widehat\mu)$, and the
participation check $v\widehat\beta$ vs $\kappa_1$.
\todo{Section~\ref{sec:empirics}: add these estimators alongside
$\widehat g$, $\widehat\alpha$.} With multiple beneficiaries,
contamination is a public good among low-floor solvers and free-riding
lowers the aggregate below the joint optimum; the largest-$v$ attacker's
bias lower-bounds the Nash aggregate. Scope: linearized sensitivity,
memoryless window, single pair, mechanical benchmark.
\end{remark}

\begin{remark}
The proposition is stated for deterministic mechanisms; whether
randomization helps (in expectation, against an oblivious adversary) is
part of Open Problem~\ref{op:game}.
\end{remark}

Consequently any deployed online notion of fairness must run against a
\emph{historical} benchmark rather than the realized one. The deployed
choice is a lower-tail quantile of past per-pair execution quality with a
\emph{downward rate limit}. Two facts about this estimator, established in
the companion empirical paper \cite{companion} and stated here for
completeness: under persistent $\alpha$-contamination of
benchmark-eligible samples, the steady-state reserve is depressed to the
$\frac{\tau-\alpha}{1-\alpha}$-quantile of the clean distribution (a bias
the bottom-dump adversary attains with equality, so the bound is tight for
this estimator); and a downward rate limit $\delta$ cannot remove
persistent bias but converts an episodic attack of window $W$ into a
bounded depression $\approx\delta W$---a \emph{persistent-bias-to-episodic-cap}
tradeoff that rationalizes the deployed design. We deliberately do not
present these as game-theoretic results: with near-costless contamination
the coalition's problem is a corner solution, and a genuine equilibrium
treatment requires a contamination-cost/detection model.

\begin{openproblem}\label{op:game}
A game-theoretic treatment of benchmark manipulation---coalition payoff
from under-delivery, equilibrium attack, and the rate limit as the
designer's best response---under an explicit contamination-cost model; and
the tight online competitive ratio for randomized online mechanisms beyond
the deterministic $g$ of Proposition~\ref{prop:online}.
\end{openproblem}

\begin{openproblem}\label{op:characterization}
A taxation-principle characterization of all feasible, fair, DSIC
mechanisms in the no-numeraire model, extending Theorem~\ref{thm:posp}
from the single-batch class to all instances
(Remark~\ref{conj:composition}); and the per-instance (rather than
worst-case) version of Theorem~\ref{thm:pof}.
\end{openproblem}

% =====================================================================
\section*{Acknowledgments}
\todo{funding; CoW practitioner discussions, if disclosable under
double-blind.}

% =====================================================================
\begin{thebibliography}{10}\small

\bibitem{canidio2024}
A.~Canidio and F.~Henneke.
\newblock Fair combinatorial auctions for trade intents.
\newblock arXiv:2408.12225, 2024.
\newblock \todo{verify exact title, author list, and latest version/venue.}

\bibitem{shapley1971}
L.~S. Shapley and M.~Shubik.
\newblock The assignment game I: the core.
\newblock {\em International Journal of Game Theory}, 1(1):111--130, 1971.

\bibitem{elyaniv2001}
R.~El-Yaniv, A.~Fiat, R.~M.~Karp, and G.~Turpin.
Optimal search and one-way trading online algorithms.
\emph{Algorithmica}, 30(1):101--139, 2001.

\bibitem{syrgkanis2013}
V.~Syrgkanis and \'E.~Tardos.
Composable and efficient mechanisms.
In \emph{Proc.\ 45th ACM STOC}, pages 211--220, 2013.

\bibitem{bikhchandani2002}
S.~Bikhchandani and J.~M. Ostroy.
\newblock The package assignment model.
\newblock {\em Journal of Economic Theory}, 107(2):377--406, 2002.

\bibitem{nisan2006}
P.~Cramton, Y.~Shoham, and R.~Steinberg, editors.
\newblock {\em Combinatorial Auctions}.
\newblock MIT Press, 2006.

\bibitem{lehmann2002}
D.~Lehmann, L.~I. O'Callaghan, and Y.~Shoham.
\newblock Truth revelation in approximately efficient combinatorial
auctions.
\newblock {\em Journal of the ACM}, 49(5):577--602, 2002.

\bibitem{karp1972}
R.~M. Karp.
\newblock Reducibility among combinatorial problems.
\newblock In {\em Complexity of Computer Computations}, pages 85--103.
Plenum, 1972.

\bibitem{bertsimas2011}
D.~Bertsimas, V.~F. Farias, and N.~Trichakis.
\newblock The price of fairness.
\newblock {\em Operations Research}, 59(1):17--31, 2011.

\bibitem{caragiannis2012}
I.~Caragiannis, C.~Kaklamanis, P.~Kanellopoulos, and M.~Kyropoulou.
\newblock The efficiency of fair division.
\newblock {\em Theory of Computing Systems}, 50(4):589--610, 2012.

\bibitem{goldberg2006}
A.~V. Goldberg, J.~D. Hartline, A.~R. Karlin, M.~Saks, and A.~Wright.
\newblock Competitive auctions.
\newblock {\em Games and Economic Behavior}, 55(2):242--269, 2006.

\bibitem{myerson1981}
R.~B. Myerson.
\newblock Optimal auction design.
\newblock {\em Mathematics of Operations Research}, 6(1):58--73, 1981.

\bibitem{companion}
Anonymous.
\newblock Solver competition under CIP-67: an empirical study.
\newblock Companion paper, anonymized for double-blind review.
\newblock \todo{cross-reference per WINE policy.}

\end{thebibliography}

\end{document}
